\documentclass{article}
\usepackage[utf8]{inputenc}
\usepackage{mathtools}
\usepackage[a4paper]{geometry}
\geometry{top=3cm, bottom=3.0cm, left=2.8cm, right=2.8cm}
\usepackage{graphicx}
\usepackage{textcomp}
\usepackage{setspace}
\usepackage{float}
\usepackage{amssymb}
\usepackage{amsmath}
\usepackage{caption}
\usepackage{subcaption}
\usepackage{pstricks}
\usepackage{fancybox}
\usepackage{comment}
\usepackage{cancel}
\usepackage{enumitem}
\usepackage{booktabs}
\usepackage{framed}

\cornersize{.2}
%\spacing{1.2}

%\newcommand{\Ns}{N_{\text{dof}}^{(\sigma)}}
\newcommand{\Ns}{N_\sigma}
\newcommand{\uh}{\mathbf{u}_h}
\newcommand{\uj}{\mathbf{u}_j}
\newcommand{\uk}{\mathbf{u}_k}
\newcommand{\sumj}{\sum_{j=0}^{\Ns}}
\newcommand{\sumk}{\sum_{k=0}^{\Ns}}
\newcommand{\nablah}{\nabla_h}
\newcommand{\pnh}{p_{\text{nh}}}
\newcommand{\pnhbottom}{p_{\text{nh},0}}
\newcommand{\ddt}[1]{\frac{\partial #1}{\partial t}}
\newcommand{\dds}[1]{\frac{\partial #1}{\partial \sigma}}
\newcommand{\phij}{\phi_j(\sigma)}
\newcommand{\phik}{\phi_k(\sigma)}
\newcommand{\phii}{\phi_i(\sigma)}
\newcommand{\vphij}{\varphi_j}
\newcommand{\vphik}{\varphi_k}
\newcommand{\Phij}{\Phi_j}
\newcommand{\Phik}{\Phi_k}
\newcommand{\Gam}{\boldsymbol{\Gamma}_{\text{mesh}}}

\newcommand{\VMij}{\mathbf{M}_{ij}^V}
\newcommand{\VAij}{\mathbf{A}_{ij}^V}
\newcommand{\VKij}{\mathbf{K}_{ij}^V}
\newcommand{\VMikj}{\boldsymbol{\mathcal{M}}_{ikj}^V}
\newcommand{\VGikj}{\boldsymbol{\mathcal{G}}_{ikj}^V}
\newcommand{\VAikj}{\boldsymbol{\mathcal{A}}_{ikj}^V}
\newcommand{\VKikj}{\boldsymbol{\mathcal{K}}_{ikj}^V}

\newcommand{\VmM}{\mathbf{M}^V}
\newcommand{\VmA}{\mathbf{A}^V}
\newcommand{\VmK}{\mathbf{K}^V}
\newcommand{\VtM}{\boldsymbol{\mathcal{M}}^V}
\newcommand{\VtG}{\boldsymbol{\mathcal{G}}^V}
\newcommand{\VtA}{\boldsymbol{\mathcal{A}}^V}
\newcommand{\VtK}{\boldsymbol{\mathcal{K}}^V}
\newcommand{\VPij}{\mathbf{P}_{ij}^V}
\newcommand{\VPikj}{\boldsymbol{\mathcal{P}}_{ikj}^V}
\newcommand{\VmP}{\mathbf{P}^V}
\newcommand{\VtP}{\boldsymbol{\mathcal{P}}^V}
\newcommand{\Pcol}{\mathcal{P}_i}

\newcommand{\vi}{\mathbf{v}_i}
\newcommand{\Oh}{\Omega_h}
\newcommand{\sumi}{\sum_{i=0}^{\Ns}}
\newcommand{\FMU}{\boldsymbol{\mathcal{F}}_{\mathcal{M}}(\mathbf{U})}
\newcommand{\FGHU}{\boldsymbol{\mathcal{F}}_{\mathcal{G}}(H,\mathbf{U})}
\newcommand{\VL}{\boldsymbol{\mathcal{L}}}
\newcommand{\VN}{\boldsymbol{\mathcal{N}}}

\newcommand{\Nh}{N_h}
\newcommand{\sumNhs}{\sum_{s=0}^{\Nh}}
\newcommand{\sumNhr}{\sum_{r=0}^{\Nh}}
\newcommand{\sumNhn}{\sum_{n=0}^{\Nh}}
\newcommand{\sumNhm}{\sum_{m=0}^{\Nh}}
\newcommand{\sumNh}[1]{\sum_{#1=0}^{\Nh}}
\newcommand{\uhjn}{\hat{\mathbf{u}}_{jn}}
\newcommand{\uhkn}{\hat{\mathbf{u}}_{kn}}
\newcommand{\uhjm}{\hat{\mathbf{u}}_{jm}}
\newcommand{\uhjr}{\hat{\mathbf{u}}_{jr}}
\newcommand{\Hm}{\hat{H}_m}
\newcommand{\Hn}{\hat{H}_n}

\newcommand{\ddxa}[1]{\frac{\partial #1}{\partial x_a}}

\newcommand{\HMsm}{\mathbf{M}_{sm}^H}
\newcommand{\HMsnm}{\boldsymbol{\mathcal{M}}_{snm}^H}
\newcommand{\HCsnm}{\boldsymbol{\mathcal{C}}_{snm}^H}
\newcommand{\HDsnm}{\boldsymbol{\mathcal{D}}_{snm}^H}
\newcommand{\HGsnmr}{\boldsymbol{\mathcal{G}_{snmr}^H}}

\newcommand{\HDsn}{\boldsymbol{\mathcal{D}}^H_{sn}}

\newcommand{\HCsnma}{\mathbb{C}_{snm}^H}
\newcommand{\HDsnma}{\mathbb{D}_{snm}^H}

\newcommand{\HmM}{\mathbf{M}^H}
\newcommand{\HtM}{\boldsymbol{\mathcal{M}}^H}
\newcommand{\HtC}{\boldsymbol{\mathcal{C}}^H}
\newcommand{\HtD}{\boldsymbol{\mathcal{D}}^H}
\newcommand{\HtG}{\boldsymbol{\mathcal{G}}^H}

\newcommand{\HtDdwH}{\bar{\boldsymbol{\mathcal{D}}}^H(\Hvec)}
\newcommand{\HDdwsn}{\bar{\boldsymbol{\mathcal{D}}}_{sn}^H}
\newcommand{\HtDdw}{\bar{\boldsymbol{\mathcal{D}}}^H}
\newcommand{\HmMdwH}{\bar{\mathbf{M}}^H(\Hvec)}

\newcommand{\HttC}{\mathbb{C}^H}
\newcommand{\HttD}{\mathbb{D}^H}

\newcommand{\Uj}{\mathbf{U}_j}
\newcommand{\Hvec}{\mathbf{H}}
\newcommand{\uvec}{\mathbf{u}_G}
\newcommand{\Ublock}{\mathbf{U}_G}

\title{New 2-D Horizontal Free-Surface-Flow Models:\\
  Detailed Derivation of the LFE-M Model}
\author{Pau Manyer Fuertes}
\date{2026}

\begin{document}

\maketitle

This document provides the step-by-step mathematical derivation of the depth-integrated, non-hydrostatic multi-layer Finite Element model (LFE-M) for a 2D horizontal layout ($\Omega_h \subset \mathbb{R}^2$), mapping directly to the framework introduced by Yang \& Liu (2024). 

\section{3D Governing Equations \& Boundary Conditions}

Let the horizontal coordinates be $\mathbf{x} = (x, y)^T$ and the vertical coordinate be $z$. 
The fluid domain is bounded above by the free surface $z = \eta(\mathbf{x}, t)$ and below by the solid bed $z = -h(\mathbf{x})$. The total water depth is $H(\mathbf{x}, t) = \eta(\mathbf{x}, t) + h(\mathbf{x})$.
The 3D velocity vector is split into a 2D horizontal velocity field $\uh = (u, v)^T$ and a scalar vertical velocity $w$.
We also define the horizontal gradient $\nablah = (\partial_x, \partial_y)^T$ in order to separate the 3D gradient into horizontal and vertical components, such that $\nabla = (\nablah, \partial_z)^T$. 

The starting point for deriving the LFE-M model are the 3D Navier-Stokes equations considering an inviscid incompressible fluid, also known as Euler equations. 
The equations split into horizontal and vertical components can be written as follows: 
\begin{subequations} \label{sys: Euler eqs}
\begin{align}
    \text{Mass Continuity:} && \nablah \cdot \uh + \frac{\partial w}{\partial z} = 0  \label{eq: Euler mass} \\[1mm]
    \text{Horizontal Momentum:} && \ddt{\uh} + (\uh \cdot \nablah)\uh + w\frac{\partial \uh}{\partial z} = -\frac{1}{\rho}\nablah p  \label{eq: Euler horizontal} \\[1mm]
    \text{Vertical Momentum:} &&\ddt{w} + \uh \cdot \nablah w + w\frac{\partial w}{\partial z} = -\frac{1}{\rho}\frac{\partial p}{\partial z} - g
    \label{eq: Euler vertical}
\end{align}
where $\rho$ and $p$ are the fluid density and pressure respectively, and $g$ is the gravitational acceleration. Kinematic boundary conditions are constrained at both sea surface ($z=\eta$) and seabed ($z=-h$):
\begin{align}
    \text{Free Surface ($z = \eta$):} && w = \ddt{\eta} + \uh \cdot \nablah \eta \label{eq: surface BC} \\
    \text{Seabed ($z = -h$):} && w = -\uh \cdot \nablah h \label{eq: seabed BC}
\end{align}
\end{subequations}

\section{The $\sigma$-Coordinate Transformation}

To map the highly irregular, time-dependent 3D fluid volume into a flat reference prism, we define the normalized vertical coordinate $\sigma \in [0, 1]$:
\begin{equation}
    \sigma \equiv \sigma(\mathbf{x},z,t) \equiv \frac{z + h(\mathbf{x})}{\eta(\mathbf{x}, t) + h(\mathbf{x})} = \frac{z + h(\mathbf{x})}{H(\mathbf{x}, t)}    
\end{equation}

Using the chain rule, the differential operators transform from $(\mathbf{x}, z, t)$ into $(\mathbf{x}, \sigma, t)$ space-time as follows:
\begin{align*}
    \frac{\partial}{\partial t}\bigg|_z &= \frac{\partial}{\partial t}\bigg|_{\sigma} - \frac{\sigma}{H}\ddt{H}\frac{\partial}{\partial \sigma} \\[1mm] 
    \nablah|_z = \nablah|_\sigma + (\nablah \sigma)\frac{\partial}{\partial \sigma} &= \nablah|_\sigma + \frac{1}{H}\left(\nablah h - \sigma \nablah H\right)\frac{\partial}{\partial \sigma} \\[1mm]
    \frac{\partial}{\partial z} &= \frac{1}{H}\frac{\partial}{\partial \sigma}
\end{align*}

Applying this mapping to the previous system of equations, the Euler equations can be rewritten according to the $\sigma$-coordinate transformation. 
The mass continuity \eqref{eq: Euler mass} turns into: 
$$\nablah|_z \cdot \uh + \frac{\partial w}{\partial z} = 0 \implies \nablah|_{\sigma} \cdot \uh + \frac{1}{H}\left(\nablah h - \sigma \nablah H\right) \cdot \dds{\uh} + \frac{1}{H}\dds{w} = 0$$
$$\implies H \nablah|_{\sigma} \cdot \uh + \left(\nablah h - \sigma \nablah H\right) \cdot \dds{\uh} + \dds{w} = 0$$
where we multiplied the whole equation by the water depth $H$. 

We identify the total material derivative as the transformed vertical velocity $\omega\equiv H \frac{D\sigma}{Dt}$ which represents velocity relative to the moving mesh layers instead of working with the physical vertical velocity $w$.
By definition, the relationship between the two velocities is:
\begin{equation*}
    \omega \equiv H \frac{D\sigma}{Dt} = H \left( \frac{\partial \sigma}{\partial t} + \uh \cdot \nablah|_\sigma \sigma + w \frac{\partial \sigma}{\partial z} \right)  = -\sigma \ddt{H} + \uh \cdot \nablah h - \sigma \uh \cdot \nablah H + w
\end{equation*}
\begin{equation}
    \boxed{\omega \equiv w -\sigma \ddt{H} + \uh \cdot \left(\nablah h - \sigma \nablah H \right)}
    \label{eq: def omega}
\end{equation}

Differentiating with respect to $\sigma$ this magnitude yields
$$\dds{\omega} = \dds{w} - \ddt{H} +\dds{\uh} \cdot \left(\nablah h - \sigma \nablah H\right) - \uh \cdot  \nablah H$$
$$\implies \dds{w} + \dds{\uh} \cdot \left(\nablah h - \sigma \nablah H\right) = \dds{\omega} + \ddt{H} + \uh \cdot  \nablah H$$
We can now substitute this expression inside the mass continuity equation in order to write it in terms of $\omega$ such that
$$\implies H \nablah|_{\sigma} \cdot \uh + \left(\nablah h - \sigma \nablah H\right) \cdot \dds{\uh} + \dds{w} = H \nablah|_{\sigma} \cdot \uh +\dds{\omega} + \ddt{H} + \uh \cdot  \nablah H = 0 $$
By regrouping the first and last term of the expression and dropping the $\sigma$ constant gradient to ease the notation, we finally get the $\sigma$ coordinate mass continuity equation:
\begin{equation*}
    \boxed{\ddt{H} + \nablah \cdot \left(H\uh\right) +\dds{\omega}  = 0}
\end{equation*} 

For the vertical momentum equation \eqref{eq: Euler vertical}, the total 3D pressure field $p$ is separated into its hydrostatic component $p_{\text{h}}$ and a dynamic, non-hydrostatic perturbation field $\pnh$ as follows:
\begin{equation}
    p(\mathbf{x}, z, t) = p_{\text{h}}(\mathbf{x}, z, t) + \pnh(\mathbf{x}, z, t) = \rho g (H - h - z) + \pnh(\mathbf{x}, z, t) = \rho g (\eta - z) + \pnh(\mathbf{x}, z, t)
    \label{eq: pressure split}
\end{equation}
where the hydrostatic pressure has been replaced by the weight of the fluid column from the seabed to a certain height $z$.
Differentiating the pressure splitting with respect to the physical coordinate $z$ gives:
$$\frac{\partial p}{\partial z} = -\rho g + \frac{\partial \pnh}{\partial z}$$
Plugging this into \eqref{eq: Euler vertical} cancels gravity, leaving:
$$\ddt{w}\bigg|_z + \uh \cdot \nablah|_z w + w \frac{\partial w}{\partial z} = -\frac{1}{\rho}\frac{\partial \pnh}{\partial z}$$
We now apply the chain rule operator for the vertical axis, yielding the following vertical momentum equation in $\sigma$ coordinate:
$$\ddt{w}\bigg|_{\sigma} - \frac{\sigma}{H}\ddt{H}\dds{w} + \uh \cdot \nablah|_\sigma w + \frac{1}{H}\uh \cdot \left(\nablah h - \sigma \nablah H\right)\dds{w} + \frac{w}{H} \dds{w} = -\frac{1}{\rho H}\dds{\pnh}$$
Notice that we can group the terms multiplied by $\dds{w}$ together:
$$\ddt{w}\bigg|_{\sigma} + \uh \cdot \nablah|_\sigma w + \frac{1}{H}\underbrace{\left[ w -\sigma\ddt{H} + \uh \cdot \left( \nablah h - \sigma \nablah H\right) \right]}_{\equiv\ \omega} \dds{w} = -\frac{1}{\rho H}\dds{\pnh}$$
Using the definition \eqref{eq: def omega} of the transformed vertical velocity $\omega$ and dropping the $\sigma$ constant gradient to ease the notation, this condenses beautifully into:
$$\boxed{\ddt{w} + \uh \cdot \nablah w + \frac{\omega}{H}\dds{w} = -\frac{1}{\rho H}\dds{\pnh}}$$

Finally, the same decomposition for the pressure field is applied to the horizontal momentum equation \eqref{eq: Euler horizontal}, such that:
$$\ddt{\uh}\bigg|_z + (\uh \cdot \nablah|_z)\uh + w\frac{\partial \uh}{\partial z} = -\frac{1}{\rho}\nablah|_z \big(\rho g (\eta - z) + \pnh\big)$$
and yields
$$\ddt{\uh}\bigg|_z + (\uh \cdot \nablah|_z)\uh + w\frac{\partial \uh}{\partial z} = - g\nablah|_z \eta -\frac{1}{\rho}\nablah|_z \pnh$$
We apply the $\sigma$ mapping to all operators in the equations:
\begin{multline*}
    \ddt{\uh}\bigg|_\sigma - \frac{\sigma}{H}\ddt{H}\dds{\uh}+(\uh \cdot \nablah|_\sigma)\uh + \frac{1}{H}\left[ \uh \cdot \Big(\nablah h - \sigma \nablah H\Big) \right] \dds{\uh} \\
    + \frac{w}{H}\dds{\uh} = - g\nablah|_\sigma \eta - \frac{1}{\rho}\nablah|_\sigma \pnh - \frac{1}{\rho H}\Big(\nablah h - \sigma \nablah H\Big)\dds{\pnh}
\end{multline*}
Notice that we can group the terms multiplied by $\dds{\uh}$ together:
\begin{multline*}
    \ddt{\uh}\bigg|_\sigma +(\uh \cdot \nablah|_\sigma)\uh + \frac{1}{H}\underbrace{\left[ w -\sigma\ddt{H} + \uh \cdot \left( \nablah h - \sigma \nablah H\right) \right]}_{= \,\omega} \dds{\uh} \\
     +g\nablah|_\sigma \eta = - \frac{1}{\rho}\nablah|_\sigma \pnh - \frac{1}{\rho H}\Big(\nablah h - \sigma \nablah H\Big)\dds{\pnh}
\end{multline*}
Using the definition \eqref{eq: def omega} of the transformed vertical velocity $\omega$ and dropping the $\sigma$ constant gradient to ease the notation, this condenses beautifully into:
\begin{equation*}
    \boxed{\ddt{\uh} +(\uh \cdot \nablah)\uh + \frac{\omega}{H}\dds{\uh} 
     +g\nablah \eta = - \frac{1}{\rho}\nablah \pnh - \frac{1}{\rho H}\Big(\nablah h - \sigma \nablah H\Big)\dds{\pnh}}
\end{equation*}

For the boundary conditions, we start with the free surface condition \eqref{eq: surface BC} at $z = \eta$ which represents the requirement that fluid particles at the surface stay on the surface:
$$w\big|_{z=\eta} = \ddt{\eta}\bigg|_{z} + \uh\big|_{z=\eta} \cdot \nablah|_{z} \eta$$
Applying the $\sigma$ mapping to both operators $\partial t$ and $\nablah$ in this case introduces no additional terms because $\eta = \eta(\mathbf{x}, t)$ is independent of $\sigma$.
By evaluating $\sigma$ at $z=\eta$, we deduce the boundary free surface boundary condition is constrained at $\sigma=1$ on the mapped domain. 
All in all, we can write:
$$w\big|_{\sigma=1} = \ddt{\eta}\bigg|_{\sigma} + \uh\big|_{\sigma=1} \cdot \nablah|_{\sigma} \eta = \ddt{\eta} + \uh\big|_{\sigma=1} \cdot \nablah \eta $$
where we have dropped the $\sigma$ constant gradient to ease the notation.
Because $H = \eta + h$, we can isolate $\eta$ as $\eta = H - h$. 
Substituting this into the time derivative and horizontal gradient gives:
$$\ddt{\eta} = \ddt{H} - \underbrace{\ddt{h}}_{=0} = \ddt{H}, \quad \text{and} \quad \nablah \eta = \nablah H - \nablah h$$
which can be plugged back into the boundary condition such that:
$$w\big|_{\sigma=1} = \ddt{H} + \uh\big|_{\sigma=1} \cdot \big(\nablah H - \nablah h\big)$$
If we evaluate the transformed vertical velocity $\omega$ at $\sigma = 1$, we get:
$$\omega\big|_{\sigma=1} = w\big|_{\sigma=1} - \underbrace{\left[ \ddt{H} + \uh\big|_{\sigma=1} \cdot \big(\nablah H - \nablah h\Big) \right]}_{=w|_{\sigma=1}} = 0$$
Notice that the bracketed collection is identical to our expression for $w\big|_{\sigma=1}$. 
Substituting it in causes the entire right-hand side to cancel out, which means the boundary condition at the free surface can simply be written as:
$$\boxed{\omega\big|_{\sigma=1} = 0}$$
Because $\sigma=1$ tracks the moving free surface exactly, no fluid particles can cross the top mesh layer. 
Relative to the moving grid, the vertical velocity is zero.


Similarly, we can apply the $\sigma$ mapping to the the seabed boundary condition \eqref{eq: seabed BC} without additional terms appearing because $h=h(\mathbf{x})$ is independent of the vertical axis. 
Evaluating $\sigma$ at $z=-h$ corresponds to $\sigma=0$, which means we can reformulate the seabed boundary condition in the mapped domain as 
$$w\big|_{\sigma=0} = -\uh\big|_{\sigma=0} \cdot \nablah\big|_{\sigma} h = -\uh\big|_{\sigma=0} \cdot \nablah h$$
Evaluating the transformed vertical velocity $\omega$ at $\sigma=0$ returns
$$\omega\big|_{\sigma=0} = w\big|_{\sigma=0} -  (0) \ddt{H} + \uh\big|_{\sigma=0}  \cdot \big( \nablah h - (0) \nablah H \big) = w\big|_{\sigma=0} + \underbrace{\uh\big|_{\sigma=0}  \cdot \nablah h}_{=-w|_{\sigma=0}} = 0$$
Notice that the bracketed collection is identical to minus our expression for $w\big|_{\sigma=0}$. 
Substituting it in causes the entire right-hand side to cancel out, which means the boundary condition at the seabed can simply be written as:
$$\boxed{\omega\big|_{\sigma=0} = 0}$$

All in all, the Euler system of equations can be written in the $\sigma$-coordinate transformed domain as follows:
\begin{subequations} \label{sys: Euler sigma eqs}
\begin{align}
    \ddt{H} + \nablah \cdot \left(H\uh\right) +\dds{\omega}  &= 0  \label{eq: Euler sigma mass} \\[1mm]
    \ddt{\uh} +(\uh \cdot \nablah)\uh + \frac{\omega}{H}\dds{\uh} 
     +g\nablah \eta &= - \frac{1}{\rho}\nablah \pnh - \frac{1}{\rho H}\Big(\nablah h - \sigma \nablah H\Big)\dds{\pnh}  \label{eq: Euler sigma horizontal} \\[1mm]
    \ddt{w} + \uh \cdot \nablah w + \frac{\omega}{H}\dds{w} &= -\frac{1}{\rho H}\dds{\pnh}
    \label{eq: Euler sigma vertical}
\end{align}
with boundary conditions:
\begin{align}
    \omega &= 0 \label{eq: sigma surface BC} \quad \text{at free surface }\sigma = 1 \\
    \omega &= 0 \quad \text{at seabed }\sigma = 0  \label{eq: sigma seabed BC}
\end{align}
\end{subequations}

\newpage

\section{The Vertical Layer Finite Element Approximation}
\label{sec: vertical FE approximation}

Following the multi-layer framework, we slice the vertical domain $\sigma \in [0,1]$ into $M$ elements (layers). We approximate the continuous horizontal velocity vector field $\uh(\mathbf{x}, \sigma, t)$ by interpolating it with a set of unidimensional finite element basis functions $\phi_j(\sigma)$ of arbitrary order $p$ across the vertical axis:
\begin{equation}
    \uh(\mathbf{x}, \sigma, t) \equiv \sumj \uj(\mathbf{x}, t) \phi_j(\sigma) = \sumj \begin{pmatrix} u_j(\mathbf{x}, t) \\ v_j(\mathbf{x}, t) \end{pmatrix} \phi_j(\sigma)
    \label{eq: uh FE vertical approximation}
\end{equation}

Here, $\uj(\mathbf{x}, t)$ represents the unknown horizontal velocity vector field belonging to vertical degree of freedom $j=0,1,...\Ns$, with $\Ns\in\mathbb{N}$ and with a total of $\Ns+1$ degrees of freedom. 
For a 1D discretisation with $M$ vertical elements of order $p$, the total number of vertical degrees of freedom is $M\times p +1$, yielding then $\Ns = M\times p$.
We pose $j=0$ the degree of freedom corresponding to the seabed ($\sigma=0$).


\section{Analytical Derivation of Vertical Velocity ($w$)}

We can derive the FE approximation for both mapped $\omega$ and physical $w$ vertical velocities starting from the mass continuity equation \eqref{eq: Euler sigma mass} obtained in previous section:
$$\ddt{H} + \nablah \cdot \left(H\uh\right) +\dds{\omega}  = 0 \implies \dds{\omega}   = -\ddt{H} - \nablah \cdot \left(H\uh\right) $$
We now integrate vertically, starting from the seabed ($\sigma = 0$) up to our chosen height:
$$\omega - \omega\big|_{\sigma=0} = -\int_0^\sigma \left[\ddt{H} + \nablah \cdot \left(H\uh\right) \right] d\sigma'$$
the boundary term at the seabed vanishes according to \eqref{eq: sigma seabed BC}, and the vertical FE approximation of the horizontal velocity \eqref{eq: uh FE vertical approximation} is substituted in the right-hand-side of the equation, yielding:
$$\omega = -\int_0^\sigma \left[\ddt{H} + \nablah \cdot \left( H \sumj \uj \phi_j \right) \right] d\sigma'$$
Now, both $H=H(\mathbf{x}, t)$ and $\uj = \uj(\mathbf{x}, t)$ are independent of vertical coordinate $\sigma$, and can therefore be taken out of the integral. 
Similarly, the FE basis $\phi_j = \phi_j(\sigma)$ depends only on $\sigma$, and therefore is unchanged under the horizontal gradient $\nablah$ influence.
The separation of variables approach leads thus to the following expression for $\omega$:
$$\omega = -\ddt{H}\int_0^\sigma d\sigma' - \sumj \nablah \cdot \left( H \uj\right) \int_0^\sigma \phi_j(\sigma')  d\sigma' $$
\begin{equation}
    \implies \qquad \boxed{\omega(\mathbf{x}, \sigma,t)  = - \sigma \ddt{H} - \sumj\nablah \cdot \left( H \uj\right) \, \vphij(\sigma)} 
    \label{eq: omega FE approx}
\end{equation}
where we have introduced the unit vertical basis 
\begin{subequations} \label{eq: def vertical unit basis}
\begin{equation}
    \vphij(\sigma) \equiv \int_0^\sigma \phi_j(\sigma')d\sigma'
    \label{eq: integral def vertical unit basis}
\end{equation}
In fact, this integral definition is mathematically equivalent to solving the following differential boundary value problem:
\begin{equation}
    \frac{d\vphij}{d\sigma} = \phi_j(\sigma), \quad \text{with } \vphij(0) = 0
    \label{eq: BVP def vertical unit basis}
\end{equation}
\end{subequations}

The expression for the physical vertical velocity $w$ can be recovered by substituting definition \eqref{eq: def omega} inside relation \eqref{eq: omega FE approx} obtained:
$$w -\sigma \ddt{H} + \uh \cdot \left(\nablah h - \sigma \nablah H \right) = - \sigma \ddt{H} - \sumj\nablah \cdot \left( H \uj\right) \, \vphij$$
$$\implies \quad w  = - \uh \cdot \left(\nablah h - \sigma \nablah H \right) - \sumj\nablah \cdot \left( H \uj\right) \, \vphij$$
We introduce now the vertical FE approximation of the horizontal velocity \eqref{eq: uh FE vertical approximation}, such that
$$w  = - \sumj \uj \phi_j \cdot \left(\nablah h - \sigma \nablah H \right) - \sumj\nablah \cdot \left( H \uj\right) \, \vphij$$
\begin{equation}
    \implies \quad \boxed{w(\mathbf{x}, \sigma,t)  = \sumj -\uj  \cdot \nablah h \,\phi_j + \uj\cdot \nablah H \,\sigma \phi_j -\nablah \cdot \left( H \uj\right) \, \vphij}
    \label{eq: w FE approx}
\end{equation}

As all terms sit inside the same summation, this equation can come in handy in the following derivations.


\section{Analytical Elimination of Non-Hydrostatic Pressure}

The non-hydrostatic pressure FE approximation can be derived in this case from the vertical momentum equation \eqref{eq: Euler sigma vertical} obtained in previous section:
$$\ddt{w} + \uh \cdot \nablah w + \frac{\omega}{H}\dds{w} = -\frac{1}{\rho H}\dds{\pnh}$$
Isolating for the non-hydrostatic pressure $\pnh$ yields the exact equation for our dynamic pressure gradient:
\begin{equation}
    \dds{\pnh} = -\rho H \left(\ddt{w} + \uh \cdot \nablah w + \frac{\omega}{H}\dds{w} \right) 
    \label{eq: nh pressure derivative}
\end{equation}

The three components inside the parentheses are computed using the previous expressions \eqref{eq: omega FE approx} and \eqref{eq: w FE approx} obtained respectively for $\omega$ and $w$. 
While $\omega$ appears as factor, as far as $w$ is concerned we need to calculate all its derivatives: temporal, horizontal gradient and with respect to $\sigma$. 
\begin{itemize}
    \item \textbf{Time derivative:}
    When taking the time derivative, the bathymetry is stationary ($\partial_t h = 0$), but the total depth $H$ and modal velocities $\uj$ evolve dynamically. Applying the product rule to the components yields:
    $$\ddt{w} = \sumj \left[ -\frac{\partial \uj}{\partial t} \cdot \nablah h \,\phi_j + \left(\frac{\partial \uj}{\partial t} \cdot \nablah H + \uj \cdot \nablah \ddt{H}\right)\sigma \phi_j - \nablah \cdot \left(H \frac{\partial \uj}{\partial t} + \ddt{H}\uj\right)\vphij \right]$$
    To preserve an invariant spatial formulation in the solver, the explicit terms containing $\partial_t H$ are substituted out using the depth-integrated continuity equation. 
    Hence, that means integrating over the vertical domain $\sigma \in[0,1]$ the mass continuity equation \eqref{eq: Euler sigma mass} obtained previously as follows:
    $$\int_0^1\ddt{H} + \nablah \cdot \left(H\uh\right) +\dds{\omega} \, d\sigma  = 0 $$
    $$\implies \quad \ddt{H} \big[\, \sigma \,\big]_0^1 + \int_0^1  \nablah \cdot \left(H\uh\right) \,d\sigma +\big[ \,\omega \,\big]_0^1 = 0 $$
    $$\implies \quad \ddt{H} = - \int_0^1  \nablah \cdot \left(H\uh\right) \,d\sigma $$
    where the boundary terms vanish according to \eqref{eq: sigma surface BC} and \eqref{eq: sigma seabed BC}.
    We now substitute the horizontal velocity by its FE approximation \eqref{eq: uh FE vertical approximation}, finally yielding:
    $$\ddt{H} = - \int_0^1  \nablah \cdot \left( H \sumj \uj \phi_j \right) \,d\sigma = - \sumj \nablah \cdot \left( H \uj \right) \int_0^1 \phi_j \,d\sigma$$
    \begin{equation}
        \implies \quad \boxed{\ddt{H} = - \sumj\nablah \cdot \left( H \uj\right) \, \Phi_j}
        \label{eq: depth integrated mass continuity}
    \end{equation}
    where we have introduced the depth-integrated unit vertical basis 
    \begin{equation}
    \Phi_j \equiv \vphij(1) = \int_0^1 \phi_j(\sigma) \,d\sigma\qquad.
    \label{eq: def depth-integrated unit basis}
    \end{equation}

    The depth-integrated mass continuity equation \eqref{eq: depth integrated mass continuity} allows an important and handy different expression for $\omega$, which is obtained by replacing \eqref{eq: depth integrated mass continuity} into its previous relation \eqref{eq: omega FE approx}:
    $$\omega = -\sigma\ddt{H} - \sumj\nablah \cdot \left( H \uj\right) \, \vphij  = -\sigma \left(- \sumj\nablah \cdot \left( H \uj\right) \, \Phi_j \right) - \sumj\nablah \cdot \left( H \uj\right) \, \vphij $$
    \begin{equation}
        \implies \qquad \boxed{\omega\equiv \sumj \bigg[ \nablah \cdot \left( H \uj\right) \Big( \sigma \Phi_j - \vphij\Big) \bigg]}
        \label{eq: omega full FE approx}
    \end{equation}

    After this parenthesis, we use relation \eqref{eq: depth integrated mass continuity} in the time derivative, which leads to
    \begin{multline*}
        \begin{aligned}
            \ddt{w} &= \sumj \Bigg[ -\frac{\partial \uj}{\partial t} \cdot \nablah h \,\phi_j + \left(\frac{\partial \uj}{\partial t} \cdot \nablah H + \uj \cdot \nablah \left[- \sumk\nablah \cdot \left( H \uk\right) \, \Phi_k \right]\right)\sigma \phi_j\\  
            & \qquad\qquad\qquad\qquad- \nablah \cdot \left(H \frac{\partial \uj}{\partial t} + \left[- \sumk\nablah \cdot \left( H \uk\right) \, \Phi_k\right]\uj\right)\vphij \Bigg] \\[1mm]
            &= \sumj \Bigg[ -\frac{\partial \uj}{\partial t} \cdot \nablah h \,\phi_j + \frac{\partial \uj}{\partial t} \cdot \nablah H \, \sigma \phi_j - \sumk \uj \cdot \nablah  \left(\nablah \cdot \left[ H \uk\right] \right) \, \sigma \Phi_k \phi_j\\  
            & \qquad\qquad\qquad\qquad - \nablah \cdot \left(H \frac{\partial \uj}{\partial t} \right)\vphij + \sumk\nablah\cdot \big( \nablah \cdot \left[ H \uk\right]\uj \big) \, \Phi_k \vphij \Bigg]
        \end{aligned}
    \end{multline*}
    \begin{multline}
        \implies \ddt{w}= \sumj \Bigg[ -\frac{\partial \uj}{\partial t} \cdot \nablah h \,\phi_j + \frac{\partial \uj}{\partial t} \cdot \nablah H \, \sigma \phi_j - \nablah \cdot \left(H \frac{\partial \uj}{\partial t} \right)\vphij \\[1mm]
        + \sumk \Big[ - \uj \cdot \nablah  \left(\nablah \cdot \left[ H \uk\right] \right) \, \sigma \Phi_k \phi_j +\nablah\cdot \big( \nablah \cdot \left[ H \uk\right]\uj \big) \, \Phi_k \vphij \Big] \Bigg]
        \label{eq: w time derivative expansion}
    \end{multline}
    
    
    \item \textbf{Vertical derivative:}
    Because of the separation of variables, the nodal velocities $\uj$, bathymetry $h$, and total depth $H$ are purely horizontal fields, thus meaning the vertical derivative acts exclusively on the vertical basis profiles, for which by definition we have:
    $$\frac{\partial \phi_j}{\partial \sigma} = \phi'_j(\sigma), \quad \frac{\partial (\sigma \phi_j)}{\partial \sigma} = \phi_j(\sigma) + \sigma \phi'_j(\sigma), \quad \dds{\vphij} = \phi_j(\sigma)$$
    
    Hence, differentiating now $w$ directly with respect to $\sigma$ yields:
    $$\dds{w} = \sumj \left[ -\uj \cdot \nablah h \,\phi'_j + \uj \cdot \nablah H \,(\phi_j + \sigma \phi'_j) - \nablah \cdot (H \uj)\,\phi_j \right]$$
    
    In \eqref{eq: nh pressure derivative}, this derivative is scaled by $\omega/H$. 
    Substituting then \eqref{eq: omega full FE approx} for $\omega$:
    \begin{equation*}
        \frac{\omega}{H}\dds{w} = \frac{1}{H}\left[ \sumj \bigg[ \nablah \cdot \left( H \uj\right) \Big( \sigma \Phi_j - \vphij\Big) \bigg] \right] \dds{w} 
    \end{equation*}
    and can be expanded using the derived expression for $\partial_\sigma w$ as follows:
    \begin{multline*}
        \frac{\omega}{H}\dds{w} = 
        \frac{1}{H}\left[ \sumj \bigg[ \nablah \cdot \left( H \uj\right) \Big( \sigma \Phi_j - \vphij\Big) \bigg] \right] \times \\
        \left[\sumk -\uk \cdot \nablah h \,\phi'_k 
        + \uk \cdot \nablah H \,(\phi_k + \sigma \phi'_k) - \nablah \cdot (H \uk)\,\phi_k \right] 
    \end{multline*}
    \begin{multline}
    \begin{aligned}
        \implies  \frac{\omega}{H}\dds{w} = \frac{1}{H} \sumj \sumk \Bigg[ 
        &-\nablah \cdot (H \uj) \Big(\uk \cdot \nablah h\Big) \left( \sigma \Phi_j \phi'_k - \vphij \phi'_k \right) \\ 
        &+ \nablah \cdot (H \uj) \Big(\uk \cdot \nablah H\Big) \left( \sigma \Phi_j \phi_k + \sigma^2 \Phi_j \phi'_k - \vphij \phi_k - \sigma \vphij \phi'_k \right) \\ 
        &- \nablah \cdot (H \uj) \left(\nablah \cdot (H \uk)\right) \left( \sigma \Phi_j \phi_k - \vphij \phi_k \right) 
        \Bigg]
        \end{aligned}
        \label{eq: w vertical derivative expansion}
    \end{multline}
    
    \item \textbf{Horizontal gradient:}
    Taking the horizontal gradient of $w$ requires distributing $\nablah$ across the spatial combinations using standard vector/tensor product rules. Because $\phi_j$ and $\vphij$ have zero horizontal spatial dependency, they pass outside the gradient operators, such that
    $$\nablah w = \sumj \left[ -\nablah(\uj \cdot \nablah h)\phi_j + \nablah(\uj \cdot \nablah H)\sigma \phi_j - \nablah\left(\nablah \cdot (H \uj)\right)\vphij \right]$$
    %Expanding the vector identities for each block we get:
    %\begin{align*}
        %\nablah(\uj \cdot \nablah h) &= (\nablah %\uj)^T \nablah h + (\nablah^2 h)\uj\\
        %\nablah(\uj \cdot \nablah H) &= (\nablah %\uj)^T \nablah H + (\nablah^2 H)\uj\\
        %\nablah(\nablah \cdot (H\uj)) &= \nablah %\left(H(\nablah \cdot \uj) + \uj \cdot %\nablah H\right)
    %\end{align*}
    
\end{itemize}

In \eqref{eq: nh pressure derivative}, the horizontal gradient is combined with the horizontal velocity in the form of a non-linear advection term $\uh \cdot \nablah w$. 
Thus, once again we introduce the FE approximation \eqref{eq: uh FE vertical approximation} and expand the advective term as follows:
\begin{equation}
    \uh \cdot \nablah w = \sumk\sumj \uk \cdot \Big[ -\nablah(\uj \cdot \nablah h)\phi_j\phi_k + \nablah(\uj \cdot \nablah H)\sigma \phi_j\phi_k - \nablah\left(\nablah \cdot (H \uj)\right)\vphij\phi_k \Big]
    \label{eq: nh pressure derivative advective term}
\end{equation}

All the ingredients have now been correctly expanded according to the vertical FE discretisation. 
Before moving forward with the derivation of the non-hydrostatic pressure, in order to ease the notation simplify the next steps, we shall group the linear and non-linear terms appearing in obtained expressions \eqref{eq: w time derivative expansion}, \eqref{eq: w vertical derivative expansion} and \eqref{eq: nh pressure derivative advective term}. 

For the linear terms appearing in the time derivative \eqref{eq: w time derivative expansion}, we introduce the following vectors $\theta_j(\sigma)$ and $\mathcal{L}_j(\uj, H)$ collecting the vertical profiles and their matching derivative linear component respectively, such that:
\begin{equation}
    \theta_j(\sigma) \equiv \begin{bmatrix} 
    \phi_j(\sigma) \\[1mm] 
    \sigma \phi_j(\sigma) \\[1mm] 
    \vphij(\sigma) 
    \end{bmatrix}
    \qquad \text{and} \qquad \mathcal{L}_j(\uj, H) \equiv \begin{bmatrix} 
    -\frac{\partial \uj}{\partial t} \cdot \nablah h\\[1mm]
    \frac{\partial \uj}{\partial t} \cdot \nablah H \\[1mm]
    -\nablah \cdot\left(H \frac{\partial \uj}{\partial t}\right) 
    \end{bmatrix}
    \label{eq: def nh pressure derivative linear terms vectors}
\end{equation}
so that
$$\left(-\frac{1}{\rho H}\dds{\pnh} \right)_{\text{linear}} = \left(\ddt{w}\right)_{\text{linear}} = \sumj \mathcal{L}_j(\uj, H)\cdot \theta _j(\sigma) = \sumj\mathcal{L}_j\cdot \theta _j$$

For the non-linear terms, we introduce in this case vectors $\Theta_{kj}(\sigma)$ and $\mathcal{N}_{kj}(\uk, \uj)$ collecting the crossed vertical profiles terms and their matching derivatives crossed non-linear components respectively, such that:
\begin{subequations}
\begin{equation}
    \Theta_{kj}(\sigma) \equiv 
    \begin{bmatrix} 
        \sigma \Phi_k \phi_j(\sigma)  \\[1.5mm] 
        \Phi_k \vphij(\sigma)  \\[1.5mm] 
        \phi_j(\sigma) \phi_k(\sigma) \\[1.5mm]
        \sigma \phi_j(\sigma) \phi_k(\sigma) \\[1.5mm]
        \vphij(\sigma) \phi_k(\sigma) \\[1.5mm]
        \sigma \Phi_j \phi'_k(\sigma) - \vphij(\sigma) \phi'_k(\sigma) \\[1.5mm]
        \sigma \Phi_j \phi_k(\sigma) + \sigma^2 \Phi_j \phi'_k(\sigma) - \vphij(\sigma) \phi_k(\sigma) - \sigma \vphij(\sigma) \phi'_k(\sigma) \\[1.5mm]
        \sigma \Phi_j \phi_k(\sigma) - \vphij(\sigma) \phi_k(\sigma) \\[1.5mm]
    \end{bmatrix}
\end{equation}
and 
\begin{equation}
    \mathcal{N}_{kj}(\uk, \uj, H) \equiv 
    \begin{bmatrix}
        -\uj \cdot \nablah \left(\nablah \cdot [H \uk]\right) \\[1.5mm]
        \nablah\cdot \big( \left(\nablah \cdot \left[ H \uk\right]\right)\uj \big) \\[1.5mm]
        -\uk \cdot \nablah(\uj \cdot \nablah h) \\[1.5mm]
        \uk \cdot \nablah(\uj \cdot \nablah H) \\[1.5mm]
        -\uk \cdot \nablah\left(\nablah \cdot [H \uj]\right) \\[1.5mm]
        -\frac{1}{H} \Big[ \nablah \cdot (H \uj) \Big] \Big[ \uk \cdot \nablah h \Big] \\[1.5mm]
        \frac{1}{H} \Big[ \nablah \cdot (H \uj) \Big] \Big[ \uk \cdot \nablah H \Big] \\[1.5mm]
        -\frac{1}{H} \Big[ \nablah \cdot (H \uj) \Big] \Big[ \nablah \cdot (H \uk) \Big] \\[1.5mm]
    \end{bmatrix}    
\end{equation}
\label{eq: def nh pressure derivative nonlinear terms vectors}
\end{subequations}
and thus we can write 
$$\left(-\frac{1}{\rho H}\dds{\pnh} \right)_{\text{nonlinear}} = \sumk \sumj \mathcal{N}_{kj}(\uk, \uj, H) \cdot \Theta_{kj}(\sigma) = \sumk \sumj \mathcal{N}_{kj}\cdot \Theta_{kj} $$
By grouping the derivations into these specific operator packages, the final expression for the vertical gradient of non-hydrostatic pressure simplifies to:
\begin{equation}
    \dds{\pnh} = -\rho H \left( \sumj \mathcal{L}_j(\uj, H) \cdot \theta_j(\sigma)  + \sumk \sumj \mathcal{N}_{kj}(\uk, \uj, H) \cdot \Theta_{kj}(\sigma) \right)    
    \label{eq: simplified nh pressure derivative}
\end{equation}
On the one hand, $\mathcal{L}_j$ represents the linear terms depending purely on local layer accelerations; on the other hand $\mathcal{N}_{kj}$ represents the non-linear advective transport tensor containing all quadratic combinations.

Notice that the separation of variables is again clearly visible in the construction of the different tensors: while $\mathcal{L}_j$ and $\mathcal{N}_{kj}$ carry the horizontal dependence, the vertical $\sigma$ coordinate is only present in $\theta_j$ and $\Theta_{kj}$. \\

With this compact notation in place, the non-hydrostatic pressure $\pnh$ can be obtained by integrating its vertical gradient, now written as \eqref{eq: simplified nh pressure derivative}, from the surface $\sigma = 1$, where the atmospheric pressure condition yields $\pnh|_{\sigma = 1}=0$, to an arbitrary depth $\sigma\in[0,1]$ as follows
$$\int_\sigma^1 \frac{\partial \pnh}{\partial \sigma'} \, d\sigma'= \pnh|_{\sigma = 1} - \pnh = -\rho H \int_\sigma^1\left( \sumj \mathcal{L}_j \cdot \theta_j  + \sumk \sumj \mathcal{N}_{kj} \cdot \Theta_{kj} \right)\, d\sigma'$$
$$\implies \quad \pnh(\mathbf{x}, \sigma, t) = \rho H \int_\sigma^1\left( \sumj \mathcal{L}_j \cdot \theta_j  + \sumk \sumj \mathcal{N}_{kj} \cdot \Theta_{kj} \right)\, d\sigma'$$
As pointed out earlier, due to the separation of variables the vertical integrals only affect $\theta_j$ and $\Theta_{kj}$. 
The rest of the factors can be moved out of them, and therefore the previous expression turns into
$$\implies \quad \pnh(\mathbf{x}, \sigma, t) = \rho H \sumj \mathcal{L}_j \cdot \int_\sigma^1\theta_j(\sigma')\, d\sigma'  + \rho H  \sumk \sumj \mathcal{N}_{kj} \cdot \int_\sigma^1\Theta_{kj}(\sigma') \, d\sigma'$$
and yields the FE vertical approximation closed-form expression for non-hydrostatic pressure without any linearization:
\begin{equation}
  \boxed{\pnh(\mathbf{x},\sigma,t) =  \rho H\sumj\mathcal{L}_j \cdot\Pi_j + \rho H\sumk\sumj \mathcal{N}_{kj}\cdot\Lambda_{kj}}
  \label{eq: pnh FE approx}
\end{equation}
where we have introduced tensors $\Pi_j$ and $\Lambda_{kj}$ defined as:
\begin{subequations}
    \begin{equation}
        \Pi_j\equiv \Pi_j(\sigma) \equiv \int_\sigma^1\theta_j(\sigma')\, d\sigma' \qquad \text{and} \qquad \Lambda_{kj} \equiv \Lambda_{kj}(\sigma) \equiv \int_\sigma^1\Theta_{kj}(\sigma') \, d\sigma'
        \label{eq: integral def vertical tensors}
    \end{equation}
    In fact, these integral definitions are respectively mathematically equivalent to solving the following differential boundary value problems:
    \begin{equation}
        \frac{d\Pi_j}{d\sigma} = -\theta_j(\sigma), \quad \text{with } \Pi_j(1) = 0 \qquad \text{and} \qquad \frac{d\Lambda_{kj}}{d\sigma} = -\Theta_{kj}(\sigma), \quad \text{with } \Lambda_{kj}(1) = 0
        \label{eq: BVP def vertical tensors}
    \end{equation}
\end{subequations}


\section{Vertical Projection and the Origin of the 2D Structural Matrices}

To eliminate the vertical dimension $\sigma$ entirely from the system, we insert the derived expressions for $\omega$ and $\pnh$ according to the introduced vertical FE discretisation, respectively \eqref{eq: omega full FE approx} and \eqref{eq: pnh FE approx},
back into the horizontal momentum equation \eqref{eq: Euler sigma horizontal}. 
Due to the structural approximations introduced in $\omega$ and $\pnh$ when discretising, the horizontal momentum is non longer conserved exactly, and thus produces a residual corresponding to the discretisation error. 
To minimise this error, the Galerkin Weighted Residual Method is applied on the system, by multiplying the entire vector equation by each vertical discretisation nodal basis function $\phi_i(\sigma)$ and integrating across the column $\sigma \in [0, 1]$.
Thus, the total number of momentum equations becomes also $\Ns+1$ (one per vertical degree of freedom $i=0,\dots,\Ns$), which makes the equation system closed.

Hence, we start off from the remaining set of equations we have left, the mapped horizontal momentum equation \eqref{eq: Euler sigma horizontal} for an incompressible, inviscid fluid:
$$\ddt{\uh} + \uh \cdot \nablah \uh + \frac{\omega}{H}\dds{\uh} + g\nablah \eta = -\frac{1}{\rho}\nablah \pnh - \frac{1}{\rho H}\dds{\pnh}\Big(\nablah h - \sigma \nablah H\Big)$$
To make this compatible with our depth-integrated matrix framework, we multiply the entire equation by the total column depth $H$:
$$H\ddt{\uh} + H\uh \cdot \nablah \uh + \omega\dds{\uh} + gH\nablah \eta = -\frac{1}{\rho}\left[ H\nablah \pnh + \dds{\pnh}\Big(\nablah h - \sigma \nablah H\Big) \right]$$

We apply the Weighted Residual Method vertically by multiplying the entire equation by a test function $\phi_i(\sigma)$ (we take it equal to the vertical basis, following the Galerkin methodology) and integrating across the vertical domain $\sigma \in [0, 1]$:
\begin{multline}
    \int_0^1 \left[ H\ddt{\uh} + H\uh \cdot \nablah \uh + \omega\dds{\uh} + gH\nablah \eta \right] \phi_i(\sigma) \, d\sigma \\
    = \int_0^1 \left[ -\frac{1}{\rho}\left( H\nablah \pnh + \dds{\pnh}\left(\nablah h - \sigma \nablah H\right) \right) \right] \phi_i(\sigma) \, d\sigma
    \label{eq: depth integrated sigma horizontal momentum}
\end{multline}



The following steps consist in taking each term appearing in the depth-integrated horizontal momentum equation \eqref{eq: depth integrated sigma horizontal momentum} and introduce the expressions obtained for the horizontal velocity $\uh$, the mapped vertical velocity $\omega$ and the non-hydrostatic pressure $\pnh$ after vertical FE discretisation, respectively \eqref{eq: uh FE vertical approximation}, \eqref{eq: omega full FE approx} and \eqref{eq: pnh FE approx}. \\

\begin{itemize}
    \item \textbf{Acceleration term:} For the left-hand side, we begin with the local acceleration term where we substitute $\uh$ for \eqref{eq: uh FE vertical approximation}:
    \begin{equation*}
        \mathbf{Acc}_i^V = \int_0^1 H \ddt{\uh} \phi_i(\sigma) \, d\sigma = \int_0^1 H \frac{\partial}{\partial t}\left( \sumj \uj \phi_j(\sigma) \right) \phi_i(\sigma) \, d\sigma = \sumj H \frac{\partial \uj}{\partial t} \underbrace{\int_0^1 \phi_i(\sigma) \phi_j(\sigma) \, d\sigma}_{\VMij} 
    \end{equation*}
    and thus yields the semi-discrete acceleration residual for the horizontal momentum equation at layer $i$:
    \begin{subequations}
        \begin{equation}
            \mathbf{Acc}_i^V \equiv \sumj H \frac{\partial \uj}{\partial t} \VMij
            \label{eq: horizontal momentum semi-discrete acceleration residual}
        \end{equation}
        where we isolated the fundamental Mass Matrix $\VmM$ with components:
    \begin{equation}
        \boxed{\VMij(\sigma) \equiv \int_0^1 \phi_i(\sigma) \phi_j(\sigma) \, d\sigma}
        \label{eq: def vertical mass matrix}
    \end{equation}
    \end{subequations}
    
    \item \textbf{Advection term:} We do the same for the advection term:
    $$\mathbf{Adv}_i^V = \int_0^1 \left[ H \uh \cdot \nablah \uh + \omega \dds{\uh} \right] \phi_i(\sigma) \, d\sigma$$
    We split this into its horizontal and transformed vertical components:
    $$\mathbf{Adv}_i^V = \underbrace{\int_0^1 H \left(\uh \cdot \nablah \uh\right) \phi_i(\sigma) \, d\sigma}_{\text{Horizontal Advection Part}} + \underbrace{\int_0^1 \omega \dds{\uh} \phi_i(\sigma) \, d\sigma}_{\text{Transformed Vertical Advection Part}}$$

    \begin{itemize}[label=$\diamond$]
        \item Substituting horizontal velocity for \eqref{eq: uh FE vertical approximation} in the horizontal advection part yields:
        \begin{align*}
            \int_0^1 H \left(\uh \cdot \nablah \uh\right) \phi_i(\sigma) \, d\sigma &= \int_0^1 H \left[\left( \sumk \uk \phi_k(\sigma) \right) \cdot \nablah \left( \sumj \uj \phi_j(\sigma) \right)\right] \phi_i(\sigma) \, d\sigma \\
            &= \sumk \sumj H \left( \uk \cdot \nablah \uj \right) \underbrace{\int_0^1 \phi_i(\sigma) \phi_k(\sigma) \phi_j(\sigma) \, d\sigma}_{\VMikj} \\
        \end{align*}
        where the nodal velocities $\uk$ and their spatial gradients $\nablah \uj$ have been pull along with the depth $H$ outside of the vertical integral since they are purely horizontal fields, isolating the static 3rd-order mass tensor $\VtM$.
        
        \item For the transformed vertical advection part, we substitute \eqref{eq: omega full FE approx} and \eqref{eq: uh FE vertical approximation} for $\omega$ and $\uh$ respectively, yielding:
        $$\int_0^1 \left( \sumk \left[ \nablah \cdot (H\uk)\Big(\sigma \Phi_k - \vphik(\sigma)\Big) \right] \right) \left( \sumj \uj \phi'_j(\sigma) \right) \phi_i(\sigma) \, d\sigma$$
        Pulling the horizontal operators outside of the vertical integrals transforms this component into:
        $$\sumk \sumj \nablah \cdot (H\uk) \uj \underbrace{ \int_0^1 \Big(\sigma \Phi_k - \vphik(\sigma)\Big)\phi'_j(\sigma) \phi_i(\sigma) \, d\sigma}_{\VGikj} = \sumk \sumj \nablah \cdot (H\uk) \uj  \, \VGikj $$
    \end{itemize}

    Combining these terms with their horizontal part yields the semi-discrete advection residual for the horizontal momentum equation at layer $i$:
    \begin{subequations}
        \begin{equation}
            \mathbf{Adv}_i^V \equiv \sumk \sumj \left[ H \VMikj (\uk \cdot \nablah \uj) + (\nablah \cdot [H \uk]) \uj \,  \VGikj\right]
            \label{eq: horizontal momentum semi-discrete advection residual}
        \end{equation}
        with 
        \begin{equation}
        \boxed{\VMikj(\sigma) \equiv \int_0^1 \phi_i(\sigma) \phi_k(\sigma) \phi_j(\sigma) \, d\sigma}
        \label{eq: def vertical mass 3rd order tensor}
        \end{equation}
        and
        \begin{equation}
            \boxed{\VGikj(\sigma) \equiv \int_0^1 \Big(\sigma \Phi_k - \vphik(\sigma)\Big)\phi'_j(\sigma) \phi_i(\sigma) \, d\sigma} 
            \label{eq: def vertical gradient 3rd order tensor}
        \end{equation}
    \end{subequations}
    
    \item \textbf{Gravity term:} When evaluating the gravity term, the test function $\phi_i(\sigma)$ is integrated over the exact same vertical unit domain $[0,1]$, corresponding to the depth-integrated vertical basis profiles $\Phi_i$ defined the in \eqref{eq: def depth-integrated unit basis}. 
    This yields the semi-discrete gravity residual for the horizontal momentum equation at layer $i$:
    \begin{equation}
        \mathbf{Grav}_i^V \equiv \int_0^1 gH\nablah\eta\,\phi_i(\sigma)\,d\sigma = gH\nablah\eta \int_0^1 \phi_i(\sigma)\,d\sigma = g H \nablah \eta\, \Phi_i
        \label{eq: horizontal momentum semi-discrete gravity residual}
    \end{equation}

    \item \textbf{Pressure terms:} For the right-hand side, the non-hydrostatic pressure horizontal and vertical gradient depth-integrated terms are separated and treated as follows:
    \begin{align*}
        \mathbf{RHS}_i^V &= -\frac{1}{\rho} \int_0^1 \left( H\nablah \pnh + \dds{\pnh}\left(\nablah h - \sigma \nablah H\right) \right)  \phi_i(\sigma) \, d\sigma \\[1mm]
        &= \underbrace{-\frac{1}{\rho} \int_0^1 H (\nablah \pnh) \phi_i \, d\sigma}_{\text{Horizontal gradient}} \; \underbrace{-\frac{1}{\rho} \int_0^1 \dds{\pnh}\Big(\nablah h - \sigma \nablah H\Big) \phi_i \, d\sigma}_{\text{Vertical gradient}}
    \end{align*}
    \begin{itemize}[label=$\diamond$]
        \item For the first term, we integrate by parts according to the following differential identity:
        $$\nablah (H \pnh) = \pnh \nablah H + H \nablah \pnh \qquad \Leftrightarrow \qquad H \nablah \pnh = \nablah (H \pnh) - \pnh \nablah H $$
        Simultaneously, in $\sigma$-coordinates the vertical integration limits are fixed constants and do not depend on the horizontal position, similarly to the vertical basis. 
        Therefore, the horizontal gradient can be placed in front of the integral such that:
        $$\int_0^1 H (\nablah \pnh) \phi_i \, d\sigma = \nablah \underbrace{\left( \int_0^1 H \pnh \phi_i \, d\sigma \right) }_{\equiv\, \Pcol} - \int_0^1 \pnh \nablah (H \phi_i) \, d\sigma$$
        The first term on the right-hand side is the horizontal gradient of the column-integrated pressure moment $\Pcol \equiv \int_0^1 H \pnh \phi_i \, d\sigma$, a purely horizontal scalar field.
        Now, shortly after the Weighted Residual Method is applied over the vertical axis (current procedure), the same method will be used in the horizontal domain to minimise the residual over the plane, therefore producing the final and full-domain problem weak formulation.
        When that happens, this gradient term is integrated by parts against the horizontal test function $\vi$:
        $$\int_{\Oh} \nablah \Pcol \cdot \vi \, d\Oh = -\int_{\Oh} \Pcol \, (\nablah \cdot \vi) \, d\Oh + \oint_{\partial\Oh} \Pcol \, (\vi \cdot \mathbf{n}) \, d\Gamma$$
        For Galerkin with essential/solid-wall boundary conditions (or behind sponge layers), the boundary term vanishes while the volume term $-\int_{\Oh} \Pcol (\nablah\cdot\vi)\, d\Oh$ survives and must be retained.
        This term is in fact the leading non-hydrostatic pressure force: on a flat bed ($\nablah h = 0$) and for small-amplitude motion ($\nablah H = \nablah\eta = O(\eta)$) it is the only $O(\eta)$ non-hydrostatic contribution, and it carries the entire frequency-dispersion behaviour of the model.

        As far as the second term is concerned, we expand it via the spatial product rule
        $$\nablah (H \phi_i) = \phi_i \nablah H + H \nablah \phi_i = \phi_i \nablah H$$
        where $\nablah\phi_i = 0$.
        Substituting this back into our definition for Term 1 cancels the negative signs and yields:
        $$\text{Horizontal gradient} = -\frac{1}{\rho} \nablah \Pcol + \frac{1}{\rho} \int_0^1 \pnh \phi_i \nablah H \, d\sigma$$
    
        \item For the second term, we eliminate the vertical gradient by integrating by parts vertically with respect to $\sigma$. 
        By posing $\Gamma = \big(\nablah h - \sigma \nablah H\big) \phi_i(\sigma)$, the integration by parts is performed according to the differential identity:
        $$\frac{\partial}{\partial \sigma}\Big( \pnh \Gamma\Big) = \dds{\pnh}\Gamma + \pnh\frac{\partial \Gamma}{\partial \sigma} \qquad \Leftrightarrow \qquad \dds{\pnh}\Gamma = \frac{\partial}{\partial \sigma}\Big( \pnh \Gamma\Big) - \pnh\frac{\partial \Gamma}{\partial \sigma} $$
        By the fundamental theorem of calculus, the first term corresponds to a vertical boundary integral: at the free surface ($\sigma = 1$), the atmospheric boundary condition dictates that the non-hydrostatic dynamic pressure component drops to zero ($\pnh = 0$); at the seabed ($\sigma = 0$), we have $\Gamma(0)=\nablah h \,\phi_i(0) = \nablah h$ only for the bottom node of the vertical discretisation, because $\phi_i(0) = 0, \, \forall i \neq 0$ and $\phi_0(0)=1$.  
        This means in fact:
        \begin{align*}
            \text{Vertical gradient} &= - \frac{1}{\rho} \Big[\pnh \Gamma \Big]_0^1+ \frac{1}{\rho} \int_0^1 \pnh \frac{\partial\Gamma}{\partial  \sigma} d\sigma \\[1mm]
            &= \frac{1}{\rho}\nablah h \,\pnh(0) \phi_i(0)+ \frac{1}{\rho} \int_0^1 \pnh \frac{\partial}{\partial \sigma} \left[ \Big(\nablah h - \sigma \nablah H\Big) \phi_i(\sigma) \right] d\sigma \\[1mm]
            &= \frac{1}{\rho}\nablah h \,\pnhbottom \,\phi_{i,0}+ \frac{1}{\rho} \int_0^1 \pnh \frac{\partial}{\partial \sigma} \left[ \Big(\nablah h - \sigma \nablah H\Big) \phi_i(\sigma) \right] d\sigma
        \end{align*}
        with $\pnhbottom = \pnh(0)$ and 
        \begin{equation}
            \phi_{i,0} \equiv \phi_i(0) = \begin{cases}
                0 & \text{if } i\neq 0 \quad \text{(inner and surface vertical DoFs)}\\[2mm]
                1 & \text{if } i= 0 \quad \text{(bottom vertical DoF)}
            \end{cases}
            \label{eq: def phii0}
        \end{equation}
        
        Applying the product rule to the inner vertical derivative:
        $$\frac{\partial}{\partial \sigma} \left[ \Big(\nablah h - \sigma \nablah H\Big) \phi_i \right] = \Big(\nablah h - \sigma \nablah H\Big) \frac{\partial \phi_i}{\partial \sigma} + \phi_i \frac{\partial}{\partial \sigma}\Big(\nablah h - \sigma \nablah H\Big)$$
        Since the 2D spatial bathymetry $h$ and depth $H$ have no vertical variation:
        $$\frac{\partial (\nablah h)}{\partial \sigma} = 0 \quad \text{and} \quad \frac{\partial (\sigma \nablah H)}{\partial \sigma} = \nablah H$$
        Thus, the inner derivative inside the vertical gradient term simplifies to:
        $$\text{Vertical gradient} = \frac{1}{\rho}\nablah h \,\pnhbottom \,\phi_{i,0}+ \frac{1}{\rho} \int_0^1 \pnh \left[ \Big(\nablah h - \sigma \nablah H\Big) \frac{\partial \phi_i}{\partial \sigma} - \phi_i \nablah H \right] d\sigma$$
    \end{itemize}

    Hence, we add now both terms back together under a single vertical integral factorized by the dynamic pressure $\pnh$, the retained leading gradient term, and the boundary term, such that:
    \begin{multline*}
        \mathbf{RHS}_i^V = -\frac{1}{\rho}\nablah\Pcol + \frac{1}{\rho}\nablah h \,\pnhbottom \,\phi_{i,0} \\[1mm]
        + \frac{1}{\rho} \int_0^1 \pnh \cdot \left[ \underbrace{\phi_i \nablah H}_{\text{Horizontal Gradient}} + \underbrace{\Big(\nablah h - \sigma \nablah H\Big) \frac{\partial \phi_i}{\partial \sigma} - \phi_i \nablah H}_{\text{Vertical Gradient}} \right] d\sigma
    \end{multline*}
    It can be noticed that the positive $\phi_i \nablah H$ term from the horizontal derivative transformation completely balances the negative $-\phi_i \nablah H$ term generated by the vertical product rule.
    This finally simplifies the full right-hand-side expression down to:
    $$\mathbf{RHS}_i^V = -\frac{1}{\rho}\nablah\Pcol + \frac{1}{\rho}\nablah h \,\pnhbottom \,\phi_{i,0}+ \underbrace{\frac{1}{\rho} \int_0^1 \pnh\Big(\nablah h - \sigma \nablah H\Big) \frac{\partial \phi_i}{\partial \sigma} d\sigma}_{\mathbf{I}_{\mathbf{RHS}_i^V}^{(1)}}$$
    For the integral term $\mathbf{I}_{\mathbf{RHS}_i^V}^{(1)}$, we proceed by introducing the previously derived expression for the non-hydrostatic pressure profile \eqref{eq: pnh FE approx}, cancelling out the fluid density $\rho$.
    Since the horizontal operator blocks $\mathcal{L}_j$ and $\mathcal{N}_{kj}$ along with the total depth $H$ are purely 2D fields, we factor them completely outside of the vertical $\sigma$-integrals:
    \begin{align*}
        \mathbf{I}_{\mathbf{RHS}_i^V}^{(1)} &=  \frac{1}{\rho} \int_0^1 \left[\left(\rho H \sumj \mathcal{L}_j\cdot\Pi_j + \rho H \sumk \sumj
        \mathcal{N}_{kj}\cdot\Lambda_{kj}\right)  \Big(\nablah h - \sigma \nablah H\Big) \frac{\partial \phi_i}{\partial \sigma} \right]d\sigma\\
        &= H \sumj \mathcal{L}_j \cdot \int_0^1 \Pi_j \frac{\partial \phi_i}{\partial \sigma} \Gam \, d\sigma +H \sumk \sumj \mathcal{N}_{kj} \cdot \int_0^1 \Lambda_{kj} \frac{\partial \phi_i}{\partial \sigma} \Gam\, d\sigma
    \end{align*}
    where we have introduced $\Gam=\nablah h - \sigma \nablah H$. 
    To identify the structural matrices, we expand the integration by parts on both vertical profiles $\Pi_j$ and $\Lambda_{kj}$:
    \begin{align*}
        \int_0^1 \Pi_j(\sigma) \frac{\partial \phi_i}{\partial \sigma} \Gam \, d\sigma &= \Big[ \Pi_j(\sigma) \phi_i(\sigma) \Gam \Big]_0^1 - \int_0^1 \frac{\partial}{\partial \sigma}\Big(\Pi_j(\sigma) \Gam\Big)\phi_i(\sigma)\,d\sigma \\
        &= - \Pi_j(0)\phi_{i,0}\nablah h + \int_0^1  \big(\Gam\theta_j + \Pi_j\nablah H\big)\phi_i\,d\sigma 
    \end{align*}
    and
    \begin{align*}
        \int_0^1 \Lambda_{kj}(\sigma) \frac{\partial \phi_i}{\partial \sigma} \Gam \, d\sigma &= \Big[ \Lambda_{kj}(\sigma) \phi_i(\sigma) \Gam \Big]_0^1 - \int_0^1  \frac{\partial}{\partial \sigma}\Big(\Lambda_{kj}(\sigma) \Gam\Big)\phi_i(\sigma)\,d\sigma \\
        &= - \Lambda_{kj}(0)\phi_{i,0}\nablah h + \int_0^1 \big(\Gam\Theta_{kj} + \Lambda_{kj}\nablah H\big)\phi_i\,d\sigma
    \end{align*}
    where we applied the product rule to the interior derivatives:
    $$\frac{\partial}{\partial \sigma}\Big(\Pi_j \Gam\Big) = \frac{\partial \Pi_j}{\partial \sigma}\Gam + \Pi_j \frac{\partial \Gam}{\partial \sigma} = -\theta_j(\sigma)\Gam - \Pi_j(\sigma)\nablah H$$
    and
    $$\frac{\partial}{\partial \sigma}\Big(\Lambda_{kj} \Gam\Big) = \frac{\partial \Lambda_{kj}}{\partial \sigma}\Gam + \Lambda_{kj} \frac{\partial \Gam}{\partial \sigma} = -\Theta_{kj}(\sigma)\Gam - \Lambda_{kj}(\sigma)\nablah H$$
    
    Substituting these back into the integrals cancels the negative signs out front, and yields for the right-hand-side integral term:
    \begin{align*}
        \mathbf{I}_{\mathbf{RHS}_i^V}^{(1)}  &= H \sumj \mathcal{L}_j \cdot \left[- \Pi_j(0)\phi_{i,0}\nablah h + \int_0^1  \big(\theta_j\Gam + \Pi_j\nablah H\big)\phi_i\,d\sigma \right] \\
        &\qquad+ H \sumk \sumj \mathcal{N}_{kj} \cdot \left[ - \Lambda_{kj}(0)\phi_{i,0}\nablah h + \int_0^1 \big(\Theta_{kj}\Gam + \Lambda_{kj}\nablah H\big)\phi_i\,d\sigma\right]  \\
        &= - \underbrace{H \left[\sumj \mathcal{L}_j \cdot  \Pi_j(0) + \sumk \sumj \mathcal{N}_{kj} \cdot \Lambda_{kj}(0) \right]}_{= \,\pnhbottom / \rho}\phi_{i,0}\nablah h + \mathbf{I}_{\mathbf{RHS}_i^V}^{(2)}
    \end{align*}
    where the first term is identified back to expression \eqref{eq: pnh FE approx} evaluated at ($\sigma=0$) divided by $\rho$, which in fact corresponds to the non-hydrostatic pressure at the seabed $\pnhbottom$.  
    This cancels out our previous boundary term in the right-hand side $\mathbf{RHS}_i^V$ such that
    $$\mathbf{RHS}_i^V = -\frac{1}{\rho}\nablah\Pcol + \frac{1}{\rho}\nablah h \,\pnhbottom \,\phi_{i,0} + \mathbf{I}_{\mathbf{RHS}_i^V}^{(1)} = -\frac{1}{\rho}\nablah\Pcol + \mathbf{I}_{\mathbf{RHS}_i^V}^{(2)}$$
    where $\mathbf{I}_{\mathbf{RHS}_i^V}^{(2)}$ yields:
    \begin{align*}
        \mathbf{I}_{\mathbf{RHS}_i^V}^{(2)} &= H \sumj \mathcal{L}_j \cdot \int_0^1 \Big[\theta_j\Gam + \Pi_j\nablah H\Big]\phi_i \, d\sigma \\
        &\qquad\qquad\qquad\qquad\qquad+ H \sumk \sumj \mathcal{N}_{kj} \cdot \int_0^1 \Big[\Theta_{kj}\Gam + \Lambda_{kj}\nablah H \Big]\phi_i\, d\sigma  \\
        &= H \sumj \mathcal{L}_j \cdot \int_0^1 \Big[\theta_j\left(\nablah h - \sigma \nablah H\right) + \Pi_j\nablah H\Big]\phi_i \, d\sigma \\
        &\qquad\qquad\qquad\qquad + H \sumk \sumj \mathcal{N}_{kj} \cdot \int_0^1 \Big[\Theta_{kj}\left(\nablah h - \sigma \nablah H\right) + \Lambda_{kj}\nablah H \Big]\phi_i\, d\sigma \\
        &= H \sumj \left[ \nablah h\,\mathcal{L}_j \cdot \int_0^1 \theta_j\phi_i \, d\sigma 
        - \nablah H\, \mathcal{L}_j \cdot \int_0^1 \sigma \theta_j \phi_i \, d\sigma 
        + \nablah H\, \mathcal{L}_j \cdot \int_0^1 \Pi_j \phi_i \, d\sigma\right] \\
        & + H \sumk \sumj\left[ \nablah h\, \mathcal{N}_{kj} \cdot \int_0^1 \Theta_{kj} \phi_i\, d\sigma 
        - \nablah H\, \mathcal{N}_{kj} \cdot\int_0^1  \sigma \Theta_{kj} \phi_i\, d\sigma 
        +\nablah H\, \mathcal{N}_{kj} \cdot \int_0^1  \Lambda_{kj}\phi_i\, d\sigma  \right] \\
        &= H \sumj \left[ \nablah h\, \mathcal{L}_j \cdot \underbrace{\int_0^1 \theta_j\phi_i \, d\sigma}_{\VAij} 
        + \nablah H\, \mathcal{L}_j \cdot \underbrace{\int_0^1 \left(\Pi_j-\sigma \theta_j\right) \phi_i \, d\sigma}_{\VKij} \right] \\
        & \qquad\qquad\qquad+ H \sumk \sumj\left[ \nablah h\, \mathcal{N}_{kj} \cdot \underbrace{\int_0^1 \Theta_{kj} \phi_i\, d\sigma}_{\VAikj} 
        + \nablah H\, \mathcal{N}_{kj} \cdot \underbrace{\int_0^1  \left( \Lambda_{kj}-\sigma \Theta_{kj}\right) \phi_i\, d\sigma}_{\VKikj} \right] 
    \end{align*}
    We still owe the leading gradient term $-\frac{1}{\rho}\nablah\Pcol$. Introducing the closed-form pressure \eqref{eq: pnh FE approx} into the column-integrated pressure moment $\Pcol$ and factoring the horizontal fields out of the vertical integral:
    $$\Pcol = \int_0^1 H \pnh \phi_i \, d\sigma = \rho H^2 \left[ \sumj \mathcal{L}_j \cdot \underbrace{\int_0^1 \Pi_j \phi_i \, d\sigma}_{\VPij} + \sumk\sumj \mathcal{N}_{kj} \cdot \underbrace{\int_0^1 \Lambda_{kj} \phi_i \, d\sigma}_{\VPikj} \right]$$
    
    Hence, we finally obtain the semi-discrete right-hand side residual for the horizontal momentum equation at layer $i$:
    \begin{subequations}
        \begin{multline}
            \implies \qquad \mathbf{RHS}_i^V \equiv -\nablah\left( H^2 \sumj \mathcal{L}_j\cdot \VPij + H^2 \sumk\sumj \mathcal{N}_{kj}\cdot \VPikj \right) \\
            + H \sumj \Big[\nablah h \big(\mathcal{L}_j\cdot \VAij\big)
            + \nablah H \big(\mathcal{L}_j \cdot\VKij\big) \Big] \\
            \qquad\qquad\qquad\qquad\qquad+ H \sumk \sumj\Big[\nablah h \big(\mathcal{N}_{kj}\cdot \VAikj \big)
            + \nablah H \big(\mathcal{N}_{kj}\cdot\VKikj\big) \Big]
            \label{eq: horizontal momentum semi-discrete RHS residual}
        \end{multline}
        where we have introduced the pressure term invariant tensors $\VmP$, $\VtP$, $\VmA$, $\VmK$, $\VtA$ and $\VtK$ with components defined as: 
        \begin{equation}
        \boxed{\VPij \equiv \int_0^1 \Pi_j \phi_i \, d\sigma} \qquad \boxed{\VAij \equiv \int_0^1 \phi_i \theta_j \, d\sigma} \qquad \boxed{\VKij \equiv \int_0^1 \left(\Pi_j-\sigma \theta_j\right) \phi_i \, d\sigma} 
        \label{eq: RHS residual 2nd order invariant tensors}
        \end{equation}
        and 
        \begin{equation}
            \boxed{\VPikj \equiv \int_0^1 \Lambda_{kj} \phi_i \, d\sigma} \qquad \boxed{\VAikj \equiv \int_0^1 \Theta_{kj} \phi_i\, d\sigma} \qquad\boxed{\VKikj \equiv \int_0^1  \left( \Lambda_{kj}-\sigma \Theta_{kj}\right) \phi_i\, d\sigma}
            \label{eq: RHS residual 3rd order invariant tensors}
        \end{equation}
    \end{subequations}
\end{itemize}
\vspace{4mm}

Before proceeding further, we turn our attention towards the definitions for $\VmK$, $\VtK$, $\VmP$ and $\VtP$ which concern the depth-integration of $\Pi_j$ and $\Lambda_{kj}$, themselves defined as vertical integrals in expressions \eqref{eq: integral def vertical tensors}.
Take for instance $\VmK$, where we replace \eqref{eq: integral def vertical tensors} inside the depth integral such that
$$\VKij \equiv \int_0^1 \left(\Pi_j-\sigma \theta_j\right) \phi_i \, d\sigma = \int_0^1 \left(\int_\sigma^1\theta_j(\sigma')\, d\sigma'-\sigma \theta_j\right) \phi_i \, d\sigma $$
The double integral term corresponds thus to 
$$I = \int_0^1 \left(\int_\sigma^1 \phi_i(\sigma) \:\theta_j(\tau) \, d\tau \right) d\sigma$$
where we have changed the dummy integration variable $\sigma'$ by $\tau$ in order to clear out the notation. 

Looking at the limits of integration in this double integral, the outer integral is defined over $0 \le \sigma \le 1$, while the inner integral spans $\sigma \le \tau \le 1$. 
Combining these inequalities defines a specific triangular region in the $(\sigma, \tau)$ plane. 
This region is bounded by the diagonal line $\tau = \sigma$, the vertical line $\tau = 1$, and the vertical axis $\sigma = 0$.
Fubini's theorem allows us to switch the order of integration from $d\tau \, d\sigma$ to $d\sigma \, d\tau$ by slicing this exact same triangular domain horizontally instead of vertically. 
To cover the entire triangle, the new outer variable $\tau$ must span $0 \le \tau \le 1$. 
Now, for any fixed value of $\tau$, $\sigma$ varies between $0 \le \sigma \le \tau$.
Rewriting the double integral with these new boundaries gives:
$$I = \int_0^1 \left( \int_0^\tau \phi_i(\sigma) \; \theta_j(\tau) \, d\sigma \right) d\tau =  \int_0^1 \theta_j(\tau) \left( \int_0^\tau \phi_i(\sigma) \, d\sigma \right) d\tau = \int_0^1 \theta_j(\tau) \; \varphi_i(\tau)\,d\tau$$
where, by pulling out $\theta_j$ of the $\sigma$ integral, the vertical unit basis $\varphi_i$ from \eqref{eq: def vertical unit basis} can be introduced in the relation. 
The same process can be done for $\VtK$ regarding $\Lambda_{kj}$, yielding thus $\VmK$ and $\VtK$ the final simplified expressions:
\begin{equation}
    \boxed{\VKij \equiv \int_0^1 \theta_j \big(\varphi_i-\sigma \phi_i \big) \, d\sigma} \qquad \qquad \boxed{\VKikj \equiv \int_0^1  \Theta_{kj} \big( \varphi_i-\sigma \phi_i \big)\, d\sigma}
    \label{eq: RHS residual simplified tensors}
\end{equation}
Identically, the leading-pressure tensors $\VmP$ and $\VtP$ (which are pure upward-integral moments, without the $\sigma\theta_j$ / $\sigma\Theta_{kj}$ correction) simplify via the same Fubini exchange to:
\begin{equation}
    \boxed{\VPij \equiv \int_0^1 \theta_j \,\varphi_i \, d\sigma} \qquad \qquad \boxed{\VPikj \equiv \int_0^1 \Theta_{kj}\, \varphi_i \, d\sigma}
    \label{eq: leading pressure simplified tensors}
\end{equation}
so that no upward integral is ever formed. Note the useful relations $\VKij = \VPij - \int_0^1 \sigma\theta_j\phi_i\,d\sigma$ and, for the third component, $\VPij[3] = \int_0^1 \varphi_j\varphi_i\,d\sigma = -B_{ij}$, where $B_{ij} \equiv -\int_0^1\varphi_i\varphi_j\,d\sigma \le 0$ is the classical (flat-bed) dispersion matrix: the retained leading term is precisely the carrier of the model's frequency dispersion.


\section{Final Vertically Discretised System of Equations}

By putting back together all the expressions for the residuals \eqref{eq: horizontal momentum semi-discrete acceleration residual}, \eqref{eq: horizontal momentum semi-discrete advection residual}, \eqref{eq: horizontal momentum semi-discrete gravity residual} and \eqref{eq: horizontal momentum semi-discrete RHS residual}, the final semi-discrete horizontal residual at vertical DoF $i$ can now be expressed as decoupled horizontal operators multiplied by the invariant vertical reference grid tensors:
$$\mathbf{Acc}_i^V + \mathbf{Adv}_i^V + \mathbf{Grav}_i^V = \mathbf{RHS}_i^V$$
which, combined with the depth integrated mass continuity equation \eqref{eq: depth integrated mass continuity}, completes the final vertically discretised system of $\Ns+2$ depth-integrated equations:
\begin{framed}
\begin{subequations}
    \begin{equation}
        \text{Mass Continuity:} \qquad \qquad \ddt{H} + \sumj\nablah \cdot \left( H \uj\right) \, \Phi_j = 0 
        \label{eq: semi-discrete mass continuity residual}
    \end{equation}
    Horizontal momentum at vertical layer $i$:
    \begin{multline}
        \sumj H \frac{\partial \uj}{\partial t} \VMij + \sumk \sumj \Big[ H \VMikj (\uk \cdot \nablah \uj) + (\nablah \cdot [H \uk]) \uj \, \VGikj \Big]+ g H \nablah \eta\, \Phi_i \\
        = -\nablah\left( H^2 \sumj \mathcal{L}_j\cdot\VPij + H^2 \sumk\sumj \mathcal{N}_{kj}\cdot\VPikj \right) \\
        + H\sumj\Big[\nablah h\big(\mathcal{L}_j \cdot\VAij\big) + \nablah H \big(\mathcal{L}_j\cdot \VKij\big) \Big] + H \sumk \sumj\Big[\nablah h  \big(\mathcal{N}_{kj}\cdot \VAikj \big) + \nablah H \big(\mathcal{N}_{kj}\cdot\VKikj\big) \Big]
        \label{eq: semi-discrete horizontal momentum residual}
    \end{multline}
    \label{eq: semi-discrete residual system}
\end{subequations}
\end{framed}


The following table summarises the expressions of the tensors introduced during the vertical discretisation:
 
\begin{table}[H]
\centering
\renewcommand{\arraystretch}{1.45}
\begin{tabular}{llll}
\toprule
Tensor & Definition & Order & Role \\
\midrule
$\VMij$ & $\int_0^1\phi_i\phi_j\,d\sigma$ & $2$ & Vertical mass \\
$\Phi_j$ & $\int_0^1\phi_j\,d\sigma$ & $1$ & Depth-average weight \\
$\VMikj$ & $\int_0^1\phi_i\phi_k\phi_j\,d\sigma$ & $3$ & Vertical advection mass \\
$\VGikj$ & $\int_0^1(\sigma\Phi_k-\varphi_k)\phi'_j\phi_i\,d\sigma$ & $3$ & Transformed vert.\ advection \\
$\VPij$ & $\int_0^1\theta_j\,\varphi_i\,d\sigma$ & $2$ & Leading pressure / dispersion ($\VPij[3]=-B_{ij}$) \\
$\VPikj$ & $\int_0^1\Theta_{kj}\,\varphi_i\,d\sigma$ & $3$ & Leading pressure, non-linear part \\
$\VAij$ & $\int_0^1\phi_i\,\theta_j\,d\sigma$ & $2$ & Linear pressure (bed slope) \\
$\VKij$ & $ \int_0^1 \theta_j \big(\varphi_i-\sigma \phi_i \big) \, d\sigma$ & $2$ & Linear pressure (surface slope) \\
$\VAikj$ & $\int_0^1\Theta_{kj}\phi_i\,d\sigma$ & $3$ & Non-linear pressure \\
$\VKikj$ & $\int_0^1  \Theta_{kj} \big( \varphi_i-\sigma \phi_i \big)\, d\sigma$ & $3$ & Non-linear pressure \\
\bottomrule
\end{tabular}
\caption{List of pre-computable static tensors for the LFE-M model. All vertical tensors depend only on the vertical mesh $\{\sigma_j\}$, the shape functions $\phi_j(\sigma)$ and its corresponding unit vertical basis $\varphi_j(\sigma)$.}
\label{tab: vertical tensors}
\end{table}

\section{Gridap Implementation}

\subsection*{Multi-field block compact form}

After discretising vertically, the problem should already be ready to be solved by Gridap. 
That is, following the general methodology in FEM literature, equations \eqref{eq: semi-discrete mass continuity residual} and \eqref{eq: semi-discrete horizontal momentum residual} contain only the FE basis discretisation in the vertical direction, and therefore the model is only partially discretised. 
In order to derive the full discrete Weak Form of the problem, the unknowns need to be also approximated in the horizontal domain according to a horizontal FE basis, which means the next step would be introducing a FEM approximation for horizontal velocity $\uj$ at layer $j$ and the water depth $H$. 

However, this horizontal discretisation is actually taken care under the hood in Gridap, and the user needs only to provide the continuous Weak Form on the problem domain, containing both continuous fields: the test function field and the trial function field. 

In our case, following the discretisation introduced in section \ref{sec: vertical FE approximation}, the 3D domain on which the model is solved is discretised vertically into $M$ elements of order $p$, resulting in a total of $\Ns +1 = M\times p +1$ degrees of freedom.
Therefore, we have $\Ns+1$ horizontal layers constituting the domain, each layer containing its own unknown field $\uj$ corresponding to the velocity field horizontal component. 
Moreover, only the top layer contains also the unknown water depth $H$, associated to the variable free-surface elevation horizontal field $\eta$, resulting thus in a total of $\Ns+2 = M\times p +2$ unknown horizontal fields, each one with its own respective trial field space.
Hence, we get $\Ns+1$ test FE spaces $\mathbf{v}_j$ for each horizontal velocity unknown field $\uj$ and an additional test FE space $q$ for the water depth $H$ (or free-surface elevation $\eta$ equivalently).

We can stack up vertically all horizontal velocity unknown and corresponding trial fields in multi-fields FE spaces in order to compact the problem later on and derive the Weak Form in algebraic for (translate the sums over different indices into contractions between tensors). 
Hence, we pose:
That is, 
\begin{equation}
    \mathbf{U} \equiv \begin{bmatrix}
        \mathbf{u}_0 & \mathbf{u}_1 & \dots & \mathbf{u}_{\Ns-1} & \mathbf{u}_{\Ns} 
    \end{bmatrix}^T \qquad  \text{and} \qquad \mathbf{V} \equiv \begin{bmatrix}
        \mathbf{v}_0 & \mathbf{v}_1 & \dots & \mathbf{v}_{\Ns-1} & \mathbf{v}_{\Ns}
    \end{bmatrix}^T
    \label{eq: def U and V horizontal velocity multifield FE spaces}
\end{equation}
For convenience, we didn't include $H$ and $q$ fields to the multi-field FE spaces $\mathbf{U}$ and $\mathbf{V}$ respectively in this derivation.

This is not the case in Gridap, where all unknown and corresponding trial fields are stored similarly in   
\texttt{MultiFieldFESpace} structures, and therefore $H$ and $q$ will be placed, for instance, at the first entry of each multi-field FE space, such that
\begin{equation}
    \boldsymbol{\xi} \equiv \begin{bmatrix}
        H & \mathbf{u}_0 & \mathbf{u}_1 & \dots & \mathbf{u}_{\Ns-1} & \mathbf{u}_{\Ns}
    \end{bmatrix}^T \qquad  \text{and} \qquad \mathbf{Q} \equiv \begin{bmatrix}
        q & \mathbf{v}_0 & \mathbf{v}_1 & \dots & \mathbf{v}_{\Ns-1} & \mathbf{v}_{\Ns}
    \end{bmatrix}^T
    \label{eq: def xi and Q multifield FE spaces}
\end{equation}
and thus $\boldsymbol{\xi}$ contains all unknown fields and $\mathbf{Q}$ contains all their corresponding trial fields. \\

Now, before diving into the final semi-discrete horizontal multilayer Weak Form derivation, let's say a few comments on Gridap \texttt{MultiFieldFESpace} solver.

In order to derive the horizontally continuous Weak From in each layer, we need the Weighted Residuals Method on each one of the $\Ns+2$ equations forming the multilayer system. 
That is, we weight each semi-discrete residual equation in \eqref{eq: semi-discrete residual system} by its corresponding trial horizontal field and integrate over the whole horizontal layer $j$ domain.
This will result in $\Ns+2$ horizontal residuals which conform the multilayer model system's Weak Form and can be solved using a non-linear FEM solver. 

However, Gridap's \texttt{MultiFieldFESpace} solver requires you to return a single, combined scalar weak form expression that represents the entire system of equations simultaneously.
That is, for a global system of $\Ns+2$ residuals $R_\alpha$, with $\alpha=0,1... \Ns+1$, such that
$$\mathbf{R} = \begin{bmatrix} 
R_0 \\ 
R_2 \\ 
\vdots \\ 
R_{\Ns} \\
R_{\Ns+1}
\end{bmatrix} = \mathbf{0}$$
Gridap expects the user to supply the global virtual work scalar, which is defined as the sum of every individual equation multiplied by its respective test function:
$$\text{Global Residual} = \int_{\Oh} R_0 v_0 \, d\Oh + \int_{\Oh} R_2 v_2 \, d\Oh + \dots + \int_{\Oh} R_{\Ns+1} v_{\Ns+1} \, d\Oh = \sum_{\alpha = 0}^{\Ns+1} \int_{\Oh} R_\alpha v_\alpha \, d\Oh$$
This total sum can be mathematically written in Gridap as the vector dot product:
\begin{equation}
    \text{Global Residual} = \int_{\Oh} \mathbf{R} \cdot \mathbf{v} \, d\Oh
    \label{eq: Gridap multifield FE space global residual}
\end{equation}
Even though you feed Gridap a single scalar expression via the dot product $\mathbf{R} \cdot \mathbf{v}$, Gridap does not lose track of the individual layer equations.
Behind the scenes, Gridap takes the functional derivative of your global residual with respect to each individual test function $v_\alpha$ one at a time.


\subsection*{Weak Form system of Residuals}

\subsubsection*{Mass Continuity Residual}

Hence, we start off with the mass continuity residual \eqref{eq: semi-discrete mass continuity residual} that we multiply by the test function $q$ and integrate over the top horizontal layer domain.
This yields:
$$\int_{\Oh} q \left(\;\ddt{H} + \sumj\nablah \cdot \left( H \uj\right) \, \Phi_j \right) d\Oh = \int_{\Oh} q \;\ddt{H} d\Oh + \sumj \left[ \int_{\Oh} q \;\nablah \cdot \left( H \uj\right)\,  d\Oh \right] \Phi_j= 0$$
and we can integrate by parts the second term according to the following identity:
$$\nablah \cdot(qH\uj) = q\nablah\cdot(H\uj) + \nablah q\cdot (H\uj) \implies  q\nablah\cdot(H\uj) = \nablah \cdot(qH\uj) -  \nablah q\cdot (H\uj) $$
$$\implies \quad \int_{\Oh} q \;\nablah \cdot \left( H \uj\right)\,  d\Omega_h = \oint_{\partial\Oh} qH\uj \; d\Gamma - \int_{\Oh} \nablah q\cdot (H\uj) \; d\Oh$$
The second term corresponds to an integral over the horizontal layer boundaries which, when prescribing essential boundary conditions and solving using the Continuous Galerkin, vanishes. 
Hence, the mass continuity equation weak form for the top layer can be written finally as 
\begin{subequations}
    \begin{equation}
        \int_{\Oh} q \;\ddt{H} d\Oh - \sumj \left[\int_{\Oh} \nablah q\cdot (H\uj) \; d\Oh\right] \Phi_j= 0
        \label{eq: semi-discrete mass continuity equation weak form}
    \end{equation}
which corresponds to the following global residual contribution from the mass continuity equation:
    \begin{equation}
        R_{mass} = \int_{\Oh} q \;\ddt{H} d\Oh - \left[\int_{\Oh} \nablah q\cdot (H\mathbf{U}) \; d\Oh\right] \cdot \boldsymbol{\Phi}
        \label{eq: multifield general residual mass continuity contribution}
    \end{equation}
where we have introduced vector of stacked depth-integrated vertical basis profiles:
    \begin{equation}
        \boldsymbol{\Phi} \equiv \begin{bmatrix}
        \Phi_0 & \Phi_1 & \dots & \Phi_{\Ns}
        \end{bmatrix}^T
        \label{eq: def array Phi}
    \end{equation}
\end{subequations}

\subsubsection*{Horizontal Momentum Residual}

We proceed similarly with the horizontal momentum semi-discrete residuals for each vertical node $i$, to which corresponds the test function $\vi$.
We also leverage the compact algebraic form introduced in \eqref{eq: def U and V horizontal velocity multifield FE spaces}, on the one hand to write the sums in the form of contractions and matrix-vector products, and on the other hand to write the system of resulting residuals as a global scalar residual that is passed to Gridap, as shown in \eqref{eq: Gridap multifield FE space global residual}.
We treat each term independently in the following:
\begin{itemize}
    \item \textbf{Acceleration term:} By posing $\dot{\mathbf{U}} = \begin{bmatrix}
        \partial_t\mathbf{u}_0 & \partial_t\mathbf{u}_1 & \dots & \partial_t\mathbf{u}_{\Ns-1} & \partial_t\mathbf{u}_{\Ns}
    \end{bmatrix}$ the horizontal velocity multi-field FE space time derivative, we can write the acceleration term contribution to the global residual in the following way:
    \begin{equation}
        R_{Acc} = \sumi \left[\int_{\Oh}\left(\sumj H \frac{\partial \uj}{\partial t} \VMij \right) \vi \; d\Oh \right] = \int_{\Oh}H\left(  \VmM \dot{\mathbf{U}} \right) \cdot \mathbf{V}  \; d\Oh 
        \label{eq: multifield general residual acceleration contribution}
    \end{equation}

    \item \textbf{Advection terms:} Both terms of the advection contribute to the global residual as:
    $$R_{Adv} = \sumi \left[\int_{\Oh}\left(\sumk \sumj \Big[ H \VMikj (\uk \cdot \nablah \uj) + (\nablah \cdot [H \uk]) \uj \, \VGikj \Big] \right) \vi \; d\Oh\right]$$
    By defining two vector-valued flux field functions, $\FMU$ and $\FGHU$, whose $i$-th components represent the active $j,k$ loops:
    \begin{subequations}
        \begin{equation}
            \big[\FMU \big]_i = \sumk \sumj \VMikj (\uk \cdot \nablah \uj) = \big[\mathbf{U}^T \VtM (\cdot \nablah \mathbf{U})\big]_i
        \end{equation}
        \begin{equation}
            \big[\FGHU\big]_i = \sumk \sumj \VGikj (\nablah \cdot [H \uk]) \uj = \big[ (\nablah \cdot [H\mathbf{U}])^T\VtG\mathbf{U}\big]_i
        \end{equation}
    \end{subequations}
    the global residual advection contribution can be written as: 
    \begin{equation}
        R_{Adv} = \int_{\Oh} \Big( H \FMU + \FGHU \Big) \cdot \mathbf{V} \; d\Oh
        \label{eq: multifield general residual advection contribution}
    \end{equation}

    \item \textbf{Gravity term:} The multi-field solver global residual contribution from the gravity term takes the following form:
    \begin{equation}
        R_{Grav} = \sumi \left[\int_{\Oh} g H \nablah \eta\, \Phi_i \vi \; d\Oh \right]= \int_{\Oh} g H \nablah \eta\, \boldsymbol{\Phi} \cdot \mathbf{V} \; d\Oh
        \label{eq: multifield general residual gravity contribution}
    \end{equation}

    \item \textbf{Pressure linear terms:} For the pressure linear terms, the global residual corresponds to:
    $$R_{lin}= \sumi \left[ \int_{\Omega_h} \left(H\sumj\Big[\nablah h\big(\mathcal{L}_j \cdot\VAij\big) + \nablah H \big(\mathcal{L}_j\cdot \VKij\big) \Big] \right) \vi \;d\Oh \right]$$
    By introducing the vector $\VL$ of stacked operator $\mathcal{L}_j$, such that 
    \begin{subequations}
    \begin{equation}
        \VL \equiv \begin{bmatrix}
            \mathcal{L}_0 & \mathcal{L}_1 & \dots & \mathcal{L}_{\Ns-1} & \mathcal{L}_{\Ns}
        \end{bmatrix}^T \quad ,
    \end{equation} 
    the sum over $j$ indices of terms $\mathcal{L}_j\cdot\VAij$ and $\mathcal{L}_j\cdot\VKij$ can be written in fact as the double contractions $\VL : \VmA$ and $\VL : \VmK$ respectively, and yields then the compact of the global residual:
    $$R_{lin}= \sumi \left[ \int_{\Omega_h} \left(H\Big[\nablah h\big(\VL : \VmA\big)_i + \nablah H \big(\VL : \VmK\big)_i \Big] \right) \vi \;d\Oh \right]$$
    \begin{equation}
        \implies \qquad R_{lin}= \int_{\Omega_h} H\Big[\nablah h\big(\VL : \VmA\big) + \nablah H \big(\VL : \VmK\big) \Big] \cdot \mathbf{V} \;d\Oh 
        \label{eq: multifield general residual linear pressure contribution}
    \end{equation}
    \end{subequations}
    
    \item \textbf{Pressure non-linear terms:}  For the pressure non-linear terms, the global residual corresponds to:
    $$R_{nonlin} = \sumi \left[ \int_{\Omega_h} \left( H \sumk \sumj\Big[\nablah h  \big(\mathcal{N}_{kj}\cdot \VAikj \big) + \nablah H \big(\mathcal{N}_{kj}\cdot\VKikj\big) \Big] \right) \vi \;d\Oh \right]$$
    By introducing the matrix $\VN$ of operators $\mathcal{N}_{kj}$, such that 
    \begin{subequations}
    \begin{equation}
        \VN \equiv \begin{bmatrix}
            \mathcal{N}_{0,0} & \mathcal{N}_{0,1} & \dots & \mathcal{N}_{0,\Ns-1} & \mathcal{N}_{0,\Ns} \\[1mm]
            \mathcal{N}_{1,0} & \mathcal{N}_{1,1} & \dots & \mathcal{N}_{1,\Ns-1} & \mathcal{N}_{1,\Ns} \\[1mm]
            \vdots & \vdots & & \vdots & \vdots \\[1mm]
            \mathcal{N}_{\Ns,0} & \mathcal{N}_{\Ns,1} & \dots & \mathcal{N}_{\Ns,\Ns-1} & \mathcal{N}_{\Ns,\Ns} 
        \end{bmatrix}    
    \end{equation}
    the sum over $j$ and $k$ indices of terms $\mathcal{N}_{kj}\cdot\VAikj$ and $\mathcal{N}_{kj}\cdot\VKikj$ can be written in fact as the triple contractions $\VN \therefore \VtA$ and $\VN \therefore \VtK$ respectively, where we have introduced notation $\therefore$ for the triple contraction of indices, and yields then the compact of the global residual:
    $$R_{nonlin} = \sumi \left[ \int_{\Omega_h} H \Big[\nablah h  \big(\VN\therefore \VtA \big)_i + \nablah H \big(\VN\therefore\VtK\big)_i \Big]  \vi \;d\Oh \right]$$
    \begin{equation}
        \implies R_{nonlin} = \int_{\Omega_h}  H \Big[\nablah h  \big(\VN\therefore \VtA \big) + \nablah H \big(\VN\therefore\VtK\big) \Big]  \cdot \mathbf{V} \;d\Oh
        \label{eq: multifield general residual nonlinear pressure contribution}
    \end{equation}
    \end{subequations}

    \item \textbf{Leading pressure term:} As anticipated in Section 6, for the pressure leading term we integrate by parts horizontally the contribution in order to move the derivative onto the test functions and separate the term into volume and boundary horizontal domain integrals.
    The boundary integral vanishes under essential/solid-wall boundary conditions when applying the Galerkin method, and we are left with the following:
    \begin{align*}
        R_{P} &= -\sumi \int_{\Oh} \nablah\left( H^2 \Big[\sumj \mathcal{L}_j\cdot\VPij + \sumk\sumj \mathcal{N}_{kj}\cdot\VPikj\Big]\right)\cdot\vi\; d\Oh \\
        &= \sumi \int_{\Oh} H^2 \Big[\sumj \mathcal{L}_j\cdot\VPij + \sumk\sumj \mathcal{N}_{kj}\cdot\VPikj\Big] \big(\nablah\cdot\vi\big)\, d\Oh
    \end{align*}
    By introducing the stacked test-divergence vector $\nablah\cdot\mathbf{V} \equiv \big[\nablah\cdot\mathbf{v}_0,\ \dots,\ \nablah\cdot\mathbf{v}_{\Ns}\big]^T$, the layer sums become the same double/triple contractions as before, yielding the compact contribution:
    \begin{equation}
        R_{P} = \int_{\Oh} H^2 \Big[ \big(\VL : \VmP\big) + \big(\VN \therefore \VtP\big) \Big] \cdot \big(\nablah\cdot\mathbf{V}\big) \; d\Oh
        \label{eq: multifield general residual leading pressure contribution}
    \end{equation}
\end{itemize}

All in all, the final expression global residual of the multifield Gridap solver results from summing all the contributions from all system residuals, that is summing expressions \eqref{eq: multifield general residual mass continuity contribution}, \eqref{eq: multifield general residual acceleration contribution}, \eqref{eq: multifield general residual advection contribution}, \eqref{eq: multifield general residual gravity contribution}, \eqref{eq: multifield general residual leading pressure contribution}, \eqref{eq: multifield general residual linear pressure contribution} and \eqref{eq: multifield general residual nonlinear pressure contribution}, yielding
$$\text{Global Residual} = R_{mass} + R_{Acc} + R_{Adv} + R_{Grav} - R_{P} - R_{lin} - R_{nonlin}$$
\begin{framed}
    \begin{multline}
        \implies \text{Global Residual} = \int_{\Oh} \Bigg[ q \;\ddt{H}  - \nablah q\cdot (H\mathbf{U}) \cdot \boldsymbol{\Phi} + \bigg( H \VmM \dot{\mathbf{U}} + H \FMU + \FGHU \\[1mm]
        + g H \nablah \eta\, \boldsymbol{\Phi}  - H\bigg[\nablah h \Big(\VL : \VmA + \VN\therefore \VtA \Big) + \nablah H \Big(\VL : \VmK + \VN\therefore\VtK \Big) \bigg] \bigg) \cdot \mathbf{V} \\[1mm]
        - H^2 \Big[ \big(\VL : \VmP\big) + \big(\VN \therefore \VtP\big) \Big] \cdot \big(\nablah\cdot\mathbf{V}\big) \Bigg] d\Oh
        \label{eq: multifield general residual}
    \end{multline}
\end{framed}



\subsection*{Field layout: why two stacked velocity fields $[\,H,\ \mathsf{U}_x,\ \mathsf{U}_y\,]$}

The global residual \eqref{eq: multifield general residual} is written in terms of the abstract stacked objects $\mathbf{U}$, $\mathbf{V}$, but it does not by itself dictate \emph{how} the $2(\Ns+1)$ scalar velocity unknowns are distributed among Gridap FE fields. Three layouts are possible, and the choice has major consequences for how the residual can be written:

\begin{enumerate}
    \item \textbf{Per-component scalar fields:} $\boldsymbol{\xi} = [\,H,\ u_0, v_0,\ u_1, v_1,\ \dots\,]$ — one scalar FE space per horizontal component per vertical node, $1 + 2(\Ns+1)$ fields in total.
    \item \textbf{Per-layer vector fields:} $\boldsymbol{\xi} = [\,H,\ \mathbf{u}_0,\ \mathbf{u}_1,\ \dots\,]$ — one $\mathbb{R}^2$-valued FE space per vertical node, $1 + (\Ns+1)$ fields. This is the ``physically natural'' choice, since each $\uj$ is a genuine 2D velocity field.
    \item \textbf{Stacked component fields (RETAINED):} $\boldsymbol{\xi} = [\,H,\ \mathsf{U}_x,\ \mathsf{U}_y\,]$ — exactly \emph{three} fields, where
    $$\mathsf{U}_x \equiv \big(u_0, u_1, \dots, u_{\Ns}\big) \in \mathbb{R}^{\Ns+1}, \qquad \mathsf{U}_y \equiv \big(v_0, v_1, \dots, v_{\Ns}\big) \in \mathbb{R}^{\Ns+1}$$
    collect the $x$- (respectively $y$-) components of \emph{all} vertical layers into a single vector-valued FE field of value dimension $\Ns+1$.
\end{enumerate}

The retained layout is dictated by one structural observation: \textbf{exactly one of the two discrete indices — the vertical layer $j$ or the spatial component $\{x,y\}$ — can live inside the FE value type, and it must be the layer index.}

\paragraph{Reason 1: the layer index must be a value-type axis.}
Every vertical sum in \eqref{eq: multifield general residual} is a contraction of the \emph{layer} index with a small constant tensor: $\VmM\dot{\mathbf{U}}$ contracts $\VMij$ with $\dot{\uj}$, the advection fluxes contract $\VMikj,\VGikj$ with pairs of layers, the pressure blocks contract $\VmA,\VmK,\VmP$ with $\VL$. If the layer index is distributed over separate FE fields (layouts 1 and 2), these sums can only be evaluated by \emph{decomposing the MultiField and looping over its entries} — e.g., for the acceleration in layout 2:
$$\texttt{acc}_i \;=\; \sum_{j=0}^{\Ns} \VMij\, \partial_t \texttt{u[2+j]} \qquad \text{(a Julia loop over $\Ns+1$ MultiField entries, per test layer $i$)},$$
which produces $O(\Ns^2)$ integrand terms for the mass matrix and $O(\Ns^3)$ for the advection. In layout 3 the layer axis is part of the pointwise \emph{value} of the field, so the same operations become single native tensor operations evaluated at each quadrature point:
$$\texttt{acc} \;=\; \VmM \cdot \partial_t\mathsf{U}_x \quad (\text{one matvec on a } \mathbb{R}^{\Ns+1}\text{-valued field}), \qquad
\big[\boldsymbol{\mathcal{F}}_{\mathcal{M}}\big]^x = \VtM \boldsymbol{:}\, \big(\mathsf{U}_x \otimes \partial_x\mathsf{U}_x + \mathsf{U}_y \otimes \partial_y\mathsf{U}_x\big),$$
where $\boldsymbol{:}$ is the native double contraction of the trailing two indices of the constant third-order tensor $\VtM$ with the rank-2 field value. No loop over layers ever appears in the residual, and the number of MultiField blocks in the assembled Jacobian drops from $\big(1+2(\Ns+1)\big)^2$ (e.g.\ $19^2 = 361$ blocks for $M=8$, $p=1$) to $3^2 = 9$, independently of $\Ns$.

\paragraph{Reason 2: fusing BOTH indices into the value type blows up the tensor rank.}
One could try to go one step further and use a \emph{single} velocity field of value type $\mathbb{R}^{2\times(\Ns+1)}$ (all components, all layers). The obstruction is the gradient rank. For the stacked field $\mathsf{U}_x$,
$$\nablah \mathsf{U}_x \in \mathbb{R}^{2\times(\Ns+1)} \quad\text{(rank 2)}, \qquad \partial_x \mathsf{U}_x = \mathbf{e}_x \cdot \nablah\mathsf{U}_x \in \mathbb{R}^{\Ns+1} \quad\text{(rank 1)},$$
so every object appearing in the residual — layer vectors, their outer products $\mathsf{U}_x \otimes \partial_x \mathsf{U}_x$, the constant contractions — has rank $\le 2$ as a field value (with constant tensors of rank $\le 3$), which is precisely the regime covered by the FE toolbox's native operator algebra ($\cdot$, $\otimes$, double contractions). For a fused field $\mathsf{W} \in \mathbb{R}^{2\times(\Ns+1)}$ instead,
$$\nablah \mathsf{W} \in \mathbb{R}^{2\times 2\times(\Ns+1)} \quad\text{(rank 3)},$$
the advection outer products become rank-4 field values, and the vertical contractions would require fourth-order constant tensors or component-extraction operations at every quadrature point — none of which is natively supported. The spatial index, by contrast, has length 2 and is \emph{asymmetric} in the equations anyway (the slope prefactors $\partial_x h$ vs $\partial_y h$, the test-divergence $\nablah\cdot\mathbf{V} = \partial_x\mathsf{W}_x + \partial_y\mathsf{W}_y$ in $R_P$ and in the gravity term treat the components differently), so writing the $x$ and $y$ momentum contributions explicitly costs two lines and is \emph{not} a loop.

\paragraph{Reason 3: boundary conditions.}
Solid-wall conditions constrain a \emph{single} component on a given boundary: $v_j = 0$ for all $j$ on the $y$-walls, $u_j = 0$ for all $j$ on the $x$-walls (closed basin). In layout 3 each of these is a plain whole-field homogeneous Dirichlet condition — the zero vector of $\mathbb{R}^{\Ns+1}$ imposed on the $\mathsf{U}_y$ (resp.\ $\mathsf{U}_x$) space on the wall tags — with no component masking machinery. A fused $\mathbb{R}^{2\times(\Ns+1)}$ field would require masks selecting half of the value components per boundary tag, with different masks meeting at the corner nodes: workable, but fragile (unconstrained corner degrees of freedom are a known catastrophic-instability class for this model).

\paragraph{Efficiency.}
The three layouts assemble the \emph{same} algebraic system up to a permutation of the degrees of freedom: same number of unknowns, same couplings (the $x$ and $y$ components interact through $\nablah\cdot(H\dot{\uj})$ and the advection regardless of the layout), hence identical sparsity and identical linear-solve cost. The efficiency difference is entirely in the \emph{assembly}: layout 3 evaluates one compact, loop-free integrand (small dense tensor operations on StaticArray values of size $\Ns+1$ and $(\Ns+1)^2$), while layouts 1–2 evaluate $O(\Ns^2)$–$O(\Ns^3)$ separate integrand terms whose per-cell temporaries grow accordingly. In the reference implementation the stacked layout reproduced the per-component solver to machine precision while running the fully nonlinear benchmark several times faster with an order of magnitude fewer allocations.




\subsection*{Implementation of the pressure operators $\mathcal{L}_j$ and $\mathcal{N}_{kj}$ in Gridap}

The pressure blocks of the global residual — $R_{lin}$, $R_{nonlin}$ and the leading term $R_P$ — are where implementation mistakes concentrate: the operators $\mathcal{L}_j$ and $\mathcal{N}_{kj}$ defined in \eqref{eq: def nh pressure derivative linear terms vectors} and \eqref{eq: def nh pressure derivative nonlinear terms vectors} contain \emph{nested} horizontal operators (gradients of divergences, gradients of directional derivatives) and products of unknown fields, which cannot be typed into a Gridap integrand verbatim. This subsection records, term by term, the exact expansions and simplifications used to reduce every component to expressions Gridap can evaluate, first in multi-field vector form and then in the component-separated form actually implemented.

Three ground rules govern everything below. Throughout, indices $a,b \in \{x, y\}$ denote horizontal components, $\partial_a \equiv (\nablah)_a$, and repeated component indices are summed.

\begin{enumerate}
    \item \textbf{Only first derivatives of the unknowns are admissible.} The horizontal FE spaces are $C^0$ Lagrangian: first derivatives of $\eta$ and $\uj$ exist element-wise and converge; \emph{second} derivatives of the unknowns (e.g.\ the free-surface Hessian $\partial_a\partial_b\,\eta$, or $\partial_a\big(\nablah\cdot[H\uk]\big)$) are not admissible in an integrand. Every nested operator must be either expanded until the surviving second derivatives land on \emph{prescribed} fields, moved onto the test function by integration by parts, or replaced by a projected/auxiliary field.
    \item \textbf{The bathymetry is prescribed analytically.} $h(\mathbf{x})$ enters as an analytic function, so $\nablah h$ and the bed Hessian $\partial_a\partial_b\, h$ are exact and always admissible. Since $H = h + \eta$, every occurrence of $\nablah H$ is evaluated as $\nablah h + \nablah\eta$ (first-order in the unknown), and $\partial_a\partial_b H = \partial_a\partial_b h + \partial_a\partial_b\eta$ is admissible \emph{only} in its bathymetric part.
    \item \textbf{Product-rule expansions are exact, not approximations.} Gridap differentiates composite integrand expressions with exact pointwise chain rules at the quadrature points (e.g.\ $\nablah(H\dot\uj) = H\nablah\dot\uj + \dot\uj\otimes\nablah H$ holds identically). The reason to expand by hand is therefore not accuracy but (i) exposing which factor carries the inadmissible derivative, and (ii) reusing common sub-expressions.
\end{enumerate}

\subsubsection*{(a) Term-by-term expansion in multi-field vector form}

\paragraph{Building-block fields.} All eight $\mathcal{N}_{kj}$ components and all three $\mathcal{L}_j$ components are built from four \emph{first-order} scalar fields per layer, which are computed once per residual evaluation and reused everywhere:
\begin{equation}
    s_k \equiv \nablah\cdot(H\uk) = H\,\nablah\cdot\uk + \uk\cdot\nablah H, \qquad
    a_k \equiv \uk\cdot\nablah h, \qquad
    b_k \equiv \uk\cdot\nablah H,
    \label{eq: pressure building blocks}
\end{equation}
together with $\nablah\cdot\uk$ itself. Note $s_k = H\,\nablah\cdot\uk + b_k$ (the product rule removes the nested divergence-of-a-product and lets $b_k$ be reused). The same fields evaluated on $\dot\uj$ are written $\dot s_j$, $\dot a_j$, $\dot b_j$.

\paragraph{The linear operator $\mathcal{L}_j$.} All three components are linear in $\dot\uj$ and contain only first derivatives — \emph{no integration by parts is needed anywhere in the linear package}:
\begin{equation}
    \mathcal{L}_j^{(1)} = -\dot a_j = -\dot\uj\cdot\nablah h, \qquad
    \mathcal{L}_j^{(2)} = \dot b_j = \dot\uj\cdot\nablah H, \qquad
    \mathcal{L}_j^{(3)} = -\dot s_j = -\big(H\,\nablah\cdot\dot\uj + \mathcal{L}_j^{(2)}\big),
    \label{eq: L implementable}
\end{equation}
where the last identity ($\dot s_j = H\nablah\cdot\dot\uj + \dot b_j$) both eliminates the nested $\nablah\cdot(H\,\cdot)$ and reuses $\mathcal{L}_j^{(2)}$.

The contraction with the vertical tensors must then be organised \emph{by component}, never by stacking: the double contraction $(\VL : \VmA)_i = \sumj \mathcal{L}_j\cdot\VAij$ expands over the 3-component index $\ell$ as
\begin{equation}
    (\VL : \VmA)_i = \sum_{\ell=1}^{3}\sumj \big(\VAij\big)_\ell\, \mathcal{L}_j^{(\ell)}
    \label{eq: L component contraction}
\end{equation}
i.e.\ three constant $(\Ns{+}1)\times(\Ns{+}1)$ matrices ${\VmA}^{(\ell)}$ and three matrix--vector products.
A tempting shortcut — building $\mathcal{L}_j$ as a single 3-component object whose entries are themselves fields or stacked layer-vectors, and contracting it against the array $\VmA$ — produces rank-inconsistent values (a ``vector of vectors'') and must be avoided: \emph{peel the component index $\ell$ off at build time}, storing ${\VmA}^{(1)},{\VmA}^{(2)},{\VmA}^{(3)}$ (and likewise ${\VmK}^{(\ell)}$, ${\VmP}^{(\ell)}$) as separate constants. The identical structure serves three residual blocks at once: $R_{lin}$ contracts $\mathcal{L}^{(\ell)}$ with ${\VmA}^{(\ell)},{\VmK}^{(\ell)}$, and the leading term $R_P$ contracts the \emph{same} $\mathcal{L}^{(\ell)}$ with ${\VmP}^{(\ell)}$.

\paragraph{The nonlinear operator $\mathcal{N}_{kj}$ — classification.} Substituting the building blocks \eqref{eq: pressure building blocks}, the eight components fall into three classes:

\emph{Class I — first-order products (components 6, 7, 8): directly implementable.}
\begin{equation}
    \mathcal{N}^{(6)}_{kj} = -\frac{1}{H}\, s_j\, a_k, \qquad
    \mathcal{N}^{(7)}_{kj} = \frac{1}{H}\, s_j\, b_k, \qquad
    \mathcal{N}^{(8)}_{kj} = -\frac{1}{H}\, s_j\, s_k .
    \label{eq: N678 implementable}
\end{equation}
Every factor is a first-order building block. Moreover, since the whole nonlinear pressure enters the residual multiplied by $H$ (cf.\ \eqref{eq: multifield general residual nonlinear pressure contribution}), the $1/H$ never needs to be formed: one implements directly the products $H\mathcal{N}^{(6)}_{kj} = -a_k s_j$, $H\mathcal{N}^{(7)}_{kj} = b_k s_j$, $H\mathcal{N}^{(8)}_{kj} = -s_k s_j$ (avoiding a needless division at every quadrature point).

\emph{Class II — one admissible second derivative (component 3).} Expanding the nested directional derivative with the identity $\partial_a( u_b\, g_b) = (\partial_a u_b)\, g_b + u_b\, (\partial_a g_b)$ applied to $g = \nablah h$:
\begin{equation}
    \mathcal{N}^{(3)}_{kj} = -\uk\cdot\nablah(\uj\cdot\nablah h)
    = -\,(\uk)_a \Big[ (\partial_a (\uj)_b)\,\partial_b h \;+\; (\uj)_b\, \underbrace{\partial_a \partial_b h}_{\text{bed Hessian}} \Big].
    \label{eq: N3 implementable}
\end{equation}
The only second derivative is $\partial_a\partial_b h$ — \emph{prescribed and analytic}, hence admissible (ground rule 2). Component 3 is therefore directly implementable despite its nested appearance. Equivalently, one may differentiate the assembled field $a_j$ directly, $\mathcal{N}^{(3)}_{kj} = -\uk\cdot\nablah a_j$: Gridap's chain rule reproduces \eqref{eq: N3 implementable} exactly, with the second derivative landing on the analytic $h$ inside $a_j$.

\emph{Class III — second derivatives of the unknowns (components 1, 2, 4, 5): NOT directly implementable.}
\begin{subequations}
\begin{align}
    \mathcal{N}^{(1)}_{kj} &= -\uj\cdot\nablah s_k = -\,(\uj)_a\, \partial_a s_k ,\\
    \mathcal{N}^{(2)}_{kj} &= \nablah\cdot\big(s_k\,\uj\big) = \uj\cdot\nablah s_k + s_k\,\nablah\cdot\uj
        = -\,\mathcal{N}^{(1)}_{kj} + s_k\,\nablah\cdot\uj ,\label{eq: N2 identity}\\
    \mathcal{N}^{(4)}_{kj} &= \uk\cdot\nablah(\uj\cdot\nablah H)
        = (\uk)_a\Big[(\partial_a(\uj)_b)\,\partial_b H + (\uj)_b\,\partial_a\partial_b H\Big] ,\label{eq: N4 expansion}\\
    \mathcal{N}^{(5)}_{kj} &= -\uk\cdot\nablah s_j = -\,(\uk)_a\,\partial_a s_j .
\end{align}
\end{subequations}
The identity \eqref{eq: N2 identity} is worth noting: component 2 is \emph{minus component 1 plus a first-order remainder}, so only \emph{one} of the two gradient-of-divergence objects ($\nablah s_k$) is ever needed for the pair $(1,2)$. The offending objects are of two kinds:
\begin{itemize}
    \item $\partial_a s_k = \partial_a\big(\nablah\cdot[H\uk]\big)$: expanding, $\partial_a s_k = \partial_a H\, \nablah\cdot\uk + H\,\partial_a(\nablah\cdot\uk) + \partial_a\uk\cdot\nablah H + \uk\cdot\partial_a\nablah H$ — it contains \emph{both} second derivatives of the velocities ($\partial_a\partial_b (\uk)_b$) and, through $\nablah H$, the free-surface Hessian $\partial_a\partial_b\,\eta$. Inadmissible on $C^0$ spaces.
    \item $\partial_a\partial_b H$ in \eqref{eq: N4 expansion}: its bathymetric part $\partial_a\partial_b h$ is admissible, but the part $\partial_a\partial_b\,\eta$ is the Hessian of the FE unknown — irreducible.
\end{itemize}

\paragraph{Treatment of Class III.} The residual multiplies $\mathcal{N}_{kj}$ by \emph{two different prefactors}, $\nablah h\,(\cdot\,\VAikj)$ and $\nablah H\,(\cdot\,\VKikj)$, and the viable technique differs between the two halves:
\begin{itemize}
    \item \emph{Bed-slope ($\nablah h$) half — integrate by parts onto the test function (EXACT).} Each Class-III term has the generic shape $\int_{\Oh} \Psi\, \big(\mathbf{u}\cdot\nablah g\big)\, d\Oh$ where $g \in \{s_k,\ \uj\cdot\nablah H\}$ is a \emph{first-order} scalar field, $\mathbf{u}$ a velocity, and $\Psi$ collects the prefactor $H\,\partial_\alpha h\,(\cdots)_c$ together with the test function. The divergence identity $\nablah\cdot(g\,\Psi\,\mathbf{u}) = \Psi\,(\mathbf{u}\cdot\nablah g) + g\,\nablah\cdot(\Psi\,\mathbf{u})$ gives
    \begin{equation}
        \int_{\Oh} \Psi\,\big(\mathbf{u}\cdot\nablah g\big)\,d\Oh
        = -\int_{\Oh} g\,\big(\mathbf{u}\cdot\nablah\Psi + \Psi\,\nablah\cdot\mathbf{u}\big)\,d\Oh
        + \oint_{\partial\Oh} g\,\Psi\,(\mathbf{u}\cdot\mathbf{n})\,ds ,
        \label{eq: class III IBP}
    \end{equation}
    where the boundary term vanishes on solid walls / behind sponge layers. On the right-hand side $g$ appears \emph{undifferentiated} (first-order), and $\nablah\Psi$ differentiates only the test function, $H$ and $h$ — all admissible, the worst object being again the analytic bed Hessian inside $\nablah(\partial_\alpha h)$. This treatment is exact: no auxiliary approximation is introduced. It is available for the $\nablah h$ half because $\Psi \propto \partial_\alpha h$ is smooth and prescribed.
    \item \emph{Surface-slope ($\nablah H$) half — projection or auxiliary field.} Here $\Psi \propto \partial_\alpha H = \partial_\alpha h + \partial_\alpha\eta$, so $\nablah\Psi$ regenerates $\partial_a\partial_b\,\eta$ — the integration by parts merely moves the irreducible Hessian from one factor to another, with no cancellation. Two sound options remain: (i) an $L^2$ projection of the first-order field $g$ (or of $\nablah g$) onto the FE space before differentiating — one small linear solve per step; or (ii) an auxiliary mixed unknown $\mathbf{Q}_k \approx \nablah s_k$ defined weakly by $\int_{\Oh}\mathbf{Q}_k\cdot\boldsymbol{\psi}\,d\Oh = -\int_{\Oh} s_k\,(\nablah\cdot\boldsymbol{\psi})\,d\Oh$ (integration by parts built into the definition), which keeps the whole system first-order and automatically differentiable.
\end{itemize}
Both Class-III halves are $O(A^3)$ in wave amplitude (quadratic operators multiplied by a slope) and are only required for finite-amplitude harmonic generation or strong shoaling; they are therefore flag-gated, while Class I (and Class II on variable beds) can remain always active at negligible cost.

\subsubsection*{(b) Component-separated expressions as implemented in the solver}

With the stacked layout $[\,H,\ \mathsf{U}_x,\ \mathsf{U}_y\,]$ and test fields $[\,q,\ \mathsf{W}_x,\ \mathsf{W}_y\,]$, the layer index disappears into the value type and the building blocks \eqref{eq: pressure building blocks} become \emph{layer-vector fields} in $\mathbb{R}^{\Ns+1}$, assembled once per residual evaluation from first derivatives only:
\begin{equation}
\begin{aligned}
    \mathsf{D} &\equiv \partial_x \mathsf{U}_x + \partial_y \mathsf{U}_y  &&(\nablah\cdot\uj \ \text{stacked}),\\
    \mathsf{a} &\equiv (\partial_x h)\,\mathsf{U}_x + (\partial_y h)\,\mathsf{U}_y  &&(a_j \ \text{stacked}),\\
    \mathsf{b} &\equiv (\partial_x H)\,\mathsf{U}_x + (\partial_y H)\,\mathsf{U}_y  &&(b_j \ \text{stacked}),\qquad \partial_a H = \partial_a h + \partial_a \eta,\\
    \mathsf{S} &\equiv H\,\mathsf{D} + \mathsf{b}  &&(s_j \ \text{stacked}),\\
    \mathsf{D}_W &\equiv \partial_x \mathsf{W}_x + \partial_y \mathsf{W}_y  &&(\text{test divergence, used by } R_P \text{ and gravity}),
\end{aligned}
\label{eq: stacked building blocks}
\end{equation}
with dotted versions $\dot{\mathsf{D}},\dot{\mathsf{a}},\dot{\mathsf{b}},\dot{\mathsf{S}}$ built identically from $\partial_t\mathsf{U}_x,\partial_t\mathsf{U}_y$. Recall $\partial_a f \equiv \mathbf{e}_a\cdot\nablah f$ (the unit vector contracts the \emph{first}, spatial, index of the gradient of a stacked field).

\paragraph{Linear package.} The three components of \eqref{eq: L implementable} become three layer-vector fields,
\begin{equation}
    \mathsf{L}^{(1)} = -\dot{\mathsf{a}}, \qquad
    \mathsf{L}^{(2)} = \dot{\mathsf{b}}, \qquad
    \mathsf{L}^{(3)} = -\big(H\dot{\mathsf{D}} + \mathsf{L}^{(2)}\big) = -\dot{\mathsf{S}},
    \label{eq: stacked L}
\end{equation}
and the component contraction \eqref{eq: L component contraction} is three constant matrix--vector products,
$$\big(\VL:\VmA\big) = \sum_{\ell=1}^{3} {\VmA}^{(\ell)}\,\mathsf{L}^{(\ell)} \in \mathbb{R}^{\Ns+1}
\qquad(\text{likewise } \VL:\VmK,\ \VL:\VmP).$$
Contracting the test in, the implemented linear slope-pressure residual block reads (subtracted, as it sits on the momentum RHS)
\begin{equation}
    R_{lin} = \int_{\Oh} H\sum_{\alpha\in\{x,y\}} \Big[\,\partial_\alpha h\; \mathsf{W}_\alpha\!\cdot\!\big(\textstyle\sum_\ell {\VmA}^{(\ell)}\mathsf{L}^{(\ell)}\big)
    + \partial_\alpha H\; \mathsf{W}_\alpha\!\cdot\!\big(\textstyle\sum_\ell {\VmK}^{(\ell)}\mathsf{L}^{(\ell)}\big)\Big]\, d\Oh ,
    \label{eq: stacked Rlin}
\end{equation}
and the leading pressure reuses the \emph{same} $\mathsf{L}^{(\ell)}$ against the test divergence (like $R_{lin}$, it is subtracted in the global residual, cf.\ \eqref{eq: multifield general residual leading pressure contribution}):
\begin{equation}
    R_{P} = \int_{\Oh} H^2\, \Big(\textstyle\sum_\ell {\VmP}^{(\ell)}\mathsf{L}^{(\ell)}\Big)\cdot\mathsf{D}_W \; d\Oh .
    \label{eq: stacked RP}
\end{equation}
Every object in \eqref{eq: stacked Rlin}--\eqref{eq: stacked RP} is a scalar, an $\mathbb{R}^{\Ns+1}$ value or an $(\Ns{+}1)\times(\Ns{+}1)$ constant: rank $\le 2$ throughout, first derivatives only.

\paragraph{Nonlinear package, Class I (implemented).} In the stacked layout the $(k,j)$ matrix structure of $\mathcal{N}_{kj}$ is realised by \emph{outer products} of layer-vectors, $(\mathbf{p}\otimes\mathbf{q})_{kj} = p_k\, q_j$ — note carefully which factor carries $k$ (the $\uk$-type factor, FIRST slot) and which carries $j$ (the $\uj$-type factor, SECOND slot), because the constant tensors are stored as $\VAikj$ with index order $[i,k,j]$ and the native double contraction $\big({\VtA}^{(c)} \boldsymbol{:}\, \mathsf{N}\big)_i = \sum_{k,j} {\VAikj}^{(c)}\, \mathsf{N}_{kj}$ contracts the \emph{trailing two} indices. From \eqref{eq: N678 implementable} (already multiplied by $H$):
\begin{equation}
    H\,\mathcal{N}^{(6)} = -\,\mathsf{a}\otimes\mathsf{S}, \qquad
    H\,\mathcal{N}^{(7)} = \mathsf{b}\otimes\mathsf{S}, \qquad
    H\,\mathcal{N}^{(8)} = -\,\mathsf{S}\otimes\mathsf{S},
    \label{eq: stacked N678}
\end{equation}
(in $\mathcal{N}^{(6)}_{kj} = -\frac1H s_j a_k$ the factor $a_k$ carries $k$ and $s_j$ carries $j$, hence $\mathsf{a}$ first, $\mathsf{S}$ second; $\mathcal{N}^{(8)}$ is symmetric so the order is immaterial there). The implemented block is then
\begin{equation}
    R_{nonlin}^{(6\text{--}8)} = \int_{\Oh} \sum_{\alpha\in\{x,y\}} \sum_{c=6}^{8}
    \Big[\, \partial_\alpha h\;\, \mathsf{W}_\alpha\!\cdot\!\big({\VtA}^{(c)} \boldsymbol{:}\, H\mathcal{N}^{(c)}\big)
    + \partial_\alpha H\;\, \mathsf{W}_\alpha\!\cdot\!\big({\VtK}^{(c)} \boldsymbol{:}\, H\mathcal{N}^{(c)}\big) \Big]\, d\Oh ,
    \label{eq: stacked Rnl68}
\end{equation}
with the overall $H$ prefactor of \eqref{eq: multifield general residual nonlinear pressure contribution} already absorbed into \eqref{eq: stacked N678}. (The equivalent form $\big(\mathsf{W}_\alpha\cdot{\VtA}^{(c)}\big)\boldsymbol{:}\,H\mathcal{N}^{(c)}$, contracting the test into the \emph{first} index beforehand, is also native; either order may be used.)

\paragraph{Nonlinear package, Classes II--III (stacked forms).} For completeness, the remaining components in stacked notation — with the inadmissible objects marked:
\begin{equation}
\begin{aligned}
    \mathcal{N}^{(1)} &= -\big[\, (\partial_x\mathsf{S})\otimes\mathsf{U}_x + (\partial_y\mathsf{S})\otimes\mathsf{U}_y \,\big] &&(\partial_a\mathsf{S}\ \text{inadmissible: } \partial^2\mathbf{u},\ \partial^2\eta)\\
    \mathcal{N}^{(2)} &= -\,\mathcal{N}^{(1)} + \mathsf{S}\otimes\mathsf{D} &&(\text{reuses } \mathcal{N}^{(1)}; \ \mathsf{S}\otimes\mathsf{D} \text{ first-order})\\
    \mathcal{N}^{(3)} &= -\big[\, \mathsf{U}_x\otimes\partial_x\mathsf{a} + \mathsf{U}_y\otimes\partial_y\mathsf{a} \,\big] &&(\partial_a\mathsf{a}\ \text{admissible: Hessian lands on } h)\\
    \mathcal{N}^{(4)} &= +\big[\, \mathsf{U}_x\otimes\partial_x\mathsf{b} + \mathsf{U}_y\otimes\partial_y\mathsf{b} \,\big] &&(\partial_a\mathsf{b}\ \text{inadmissible: contains } \partial^2\eta)\\
    \mathcal{N}^{(5)} &= -\big[\, \mathsf{U}_x\otimes\partial_x\mathsf{S} + \mathsf{U}_y\otimes\partial_y\mathsf{S} \,\big] = \big(\mathcal{N}^{(1)}\big)^{\!T} &&(\partial_a\mathsf{S}\ \text{inadmissible, as in } \mathcal{N}^{(1)})
\end{aligned}
\label{eq: stacked N1to5}
\end{equation}
Watch the slot ordering: in $\mathcal{N}^{(1)}_{kj} = -(\uj)_a\,\partial_a s_k$ the differentiated divergence carries $k$ (first slot of the outer product) and the velocity carries $j$, whereas in $\mathcal{N}^{(3,4,5)}_{kj} = \mp(\uk)_a\,\partial_a(\cdot)_j$ the velocity carries $k$ — so $\mathcal{N}^{(5)}$ is exactly the \emph{transpose} of $\mathcal{N}^{(1)}$, and swapping a single outer-product order contracts the wrong layer pairs against $\VAikj$. Such a mistake is invisible in every symmetric test configuration (equal layer velocities, symmetric tensors), which is precisely why the $[i,k,j]$ ordering convention must be fixed once, written down, and verified against a deliberately \emph{asymmetric} state. These components enter the residual through the Class-III treatments of subsection (a): the $\nablah h$ half via \eqref{eq: class III IBP} (exact), the $\nablah H$ half via projection/auxiliary field, both flag-gated.

\paragraph{Summary of what each residual block consumes.}
$$
\begin{array}{lll}
\textbf{block} & \textbf{fields used} & \textbf{status}\\[1mm]
R_{lin} \ \eqref{eq: stacked Rlin} & \mathsf{L}^{(1,2,3)} \ (= -\dot{\mathsf{a}},\ \dot{\mathsf{b}},\ -\dot{\mathsf{S}}) & \text{implemented, exact, first-order}\\
R_{P} \ \eqref{eq: stacked RP} & \text{same } \mathsf{L}^{(\ell)},\ \mathsf{D}_W & \text{implemented, exact, first-order}\\
R_{nonlin},\ c=6,7,8 \ \eqref{eq: stacked Rnl68} & \mathsf{a},\mathsf{b},\mathsf{S} \text{ (outer products)} & \text{implemented, exact, first-order}\\
R_{nonlin},\ c=3 & \mathsf{U},\ \partial\mathsf{a} \ (\text{bed Hessian}) & \text{admissible, direct}\\
R_{nonlin},\ c=1,2,5 & \partial\mathsf{S} & \nablah h: \text{IBP \eqref{eq: class III IBP}};\ \nablah H: \text{projection/aux.}\\
R_{nonlin},\ c=4 & \partial\mathsf{b} \ (\supset \partial^2\eta) & \nablah h: \text{IBP};\ \nablah H: \text{projection/aux.}
\end{array}
$$





\subsection*{Unit vertical basis $\{\vphij\}$ calculation}

Trick for computing unit vertical basis functions $\vphij$: recall that the integral definition \eqref{eq: integral def vertical unit basis} can be in fact also written in differential form as the boundary value problem:
\begin{equation*}
    \frac{d\vphij}{d\sigma} = \phi_j(\sigma), \quad \text{with } \vphij(0) = 0
\end{equation*}
Hence, we solve this as a standard FEM problem using the Weighted Residuals Method using Gridap. 
By definition, if $\phij$ are degree $p$ basis functions, then $\vphij$ constitutes a basis of degree $p+1$. 
Thus we pose this vertical BVP Weak Form as
$$\int_0^1 \frac{d\xi}{d\sigma}\frac{d\vphij}{d\sigma} \;d\sigma = \int_0^1 \frac{d\xi}{d\sigma}\phi_j \;d\sigma$$
with $\xi$ the test functions.


\subsection*{Wave Generation, Absorption, and Boundary Conditions}

The global residual \eqref{eq: multifield general residual} derived above describes wave
\emph{propagation} in the interior of $\Oh$, but a usable simulation additionally requires a way to
\emph{inject} energy into the domain, to \emph{absorb} it before it reaches the artificial mesh
edges, and to prescribe what happens \emph{on} those edges. In the solver these are handled by three
mechanisms, each of which augments the weak form with one explicit, closed-form term. Crucially, all
three are \emph{linear} in the unknowns, so they leave the Newton solve structure of
this section unchanged.

\subsubsection*{Boundary conditions}

Several boundary integrals were generated and then discarded during the weak-form derivation: the
continuity flux $\oint_{\partial\Oh} q H\uj\cdot\mathbf{n}\,d\Gamma$ arising from the divergence
integration by parts \eqref{eq: semi-discrete mass continuity equation weak form}; the
leading-pressure flux
$\oint_{\partial\Oh} H^2\big[(\VL:\VmP)+(\VN\therefore\VtP)\big]\cdot(\mathbf{V}\cdot\mathbf{n})\,d\Gamma$
(with $(\mathbf{V}\cdot\mathbf{n})_i = \vi\cdot\mathbf{n}$) arising from the integration by parts of
$R_P$ \eqref{eq: multifield general residual leading pressure contribution}; and — whenever the
Class-III non-linear pressure components are active — the boundary terms of their test-function
integration by parts \eqref{eq: class III IBP}. The
solver realises exactly the conditions under which those integrals vanish:

\begin{itemize}
    \item \textbf{Solid walls (essential / Dirichlet).} On a rigid lateral wall the no-penetration
    condition $\uj\cdot\mathbf{n}=0$ holds for \emph{every} vertical layer $j$. For a wall normal to
    $y$ this means the essential constraint
    \begin{equation*}
        u_j^y\big|_{\partial\Oh^{\text{wall}}} = 0, \qquad j = 0,1,\dots,\Ns,
    \end{equation*}
    imposed directly on the corresponding components of the \texttt{MultiFieldFESpace} (and
    symmetrically $u_j^x=0$ for a wall normal to $x$ when a fully closed basin is required). Both
    boundary integrals then vanish identically, and the free-surface elevation $\eta$ (equivalently
    $H$) is left \emph{unconstrained} --- it satisfies a natural (Neumann) condition, as is physically
    correct for a moving free surface against a vertical wall.

    \textbf{Implementation note.} On a Cartesian mesh the wall is tagged by its edge \emph{and its two
    end corners}; omitting the corner tags leaves the corner degrees of freedom unconstrained and
    triggers an exponential instability. The solid-wall constraint therefore uses the full tag set
    (four corners $+$ the two edge interiors) per wall.

    \item \textbf{Open boundaries (natural).} At an open end of a flume no essential condition is
    prescribed; the boundary flux is left as the natural (homogeneous) term. Physical radiation of
    the outgoing waves through such an open edge is \emph{not} enforced by the boundary condition
    itself but by placing a sponge layer (below) in front of it, so that the wave amplitude reaching
    the edge is already negligible and the residual boundary reflection is harmless.
\end{itemize}

\subsubsection*{Internal wavemaker: a mass source in the continuity equation}

Rather than prescribing an inflow boundary condition, waves are generated \emph{internally} by adding
a localised mass source $S(\mathbf{x},t)$ to the right-hand side of the depth-integrated mass
continuity equation \eqref{eq: semi-discrete mass continuity residual}:
\begin{equation}
    \ddt{H} + \sumj\nablah\cdot\left(H\uj\right)\Phij = S(\mathbf{x},t).
    \label{eq: continuity with source}
\end{equation}
Weighting by the test function $q$ and integrating over $\Oh$ exactly as in
\eqref{eq: semi-discrete mass continuity equation weak form}, the source contributes a single extra
term to the global residual:
\begin{equation}
    \boxed{\; R_{\text{wm}} = -\int_{\Oh} q\, S(\mathbf{x},t)\, d\Oh \;}
    \label{eq: wavemaker residual contribution}
\end{equation}
The source is a time-harmonic Gaussian pulse in space,
\begin{equation}
    S(\mathbf{x},t) = 2\,A\,\omega\; \mathcal{G}(\mathbf{x})\;\cos(\omega t),
    \qquad \omega = \frac{2\pi}{T},
    \label{eq: wavemaker source profile}
\end{equation}
with target wave amplitude $A$ and angular frequency $\omega$ set by the wave period $T$. The spatial
footprint $\mathcal{G}(\mathbf{x})$ selects the type of generated wavefield:
\begin{subequations}\label{eq: wavemaker footprints}
    \begin{align}
        \text{Line source (long-crested / plane waves):} && \mathcal{G}(\mathbf{x}) &= \exp\!\left[-\left(\frac{x-x_{\text{wm}}}{\sigma_{\text{wm}}}\right)^{2}\right], \label{eq: wavemaker line}\\[1mm]
        \text{Point source (radial / ring waves):} && \mathcal{G}(\mathbf{x}) &= \exp\!\left[-\frac{(x-x_{\text{wm}})^{2}+(y-y_{\text{wm}})^{2}}{\sigma_{\text{wm}}^{2}}\right]. \label{eq: wavemaker point}
    \end{align}
\end{subequations}
Here $(x_{\text{wm}}, y_{\text{wm}})$ is the generator location, and $\sigma_{\text{wm}}$ is the
Gaussian half-width controlling how sharply the source is concentrated (a narrow footprint,
$\sigma_{\text{wm}}\sim1.5$~m by default, approximates a clean single-frequency paddle while
remaining smooth enough to be resolved by the horizontal mesh). The line source
\eqref{eq: wavemaker line} has no $y$-dependence and therefore emits straight wave crests; the point
source \eqref{eq: wavemaker point} emits circular fronts.

\paragraph{The factor of 2.} Because the source injects mass symmetrically, it radiates waves in
\emph{both} directions away from the generator line/point. Only half of the emitted energy travels
toward the region of interest; the leading factor $2$ in \eqref{eq: wavemaker source profile}
compensates for this split so that the wave reaching the working area has the prescribed amplitude
$A$. The counter-propagating half is absorbed by the sponge layer on the opposite side.

\paragraph{Solver implementation.} The profile \eqref{eq: wavemaker source profile} is built by the
factory routines \texttt{make\_wavemaker\_line} and \texttt{make\_wavemaker\_point} (arguments
$x_{\text{wm}}$, $y_{\text{wm}}$, $A$, $T$, $k$ and $\sigma_{\text{wm}}$), which return a
closure $S(\mathbf{x},t)$. This closure is stored on the problem struct (\texttt{prob.wm\_src}) and,
at each time $t$, wrapped into a Gridap \texttt{CellField} evaluated on the integration mesh; the term
\eqref{eq: wavemaker residual contribution} is then appended to the continuity residual.

\subsubsection*{Sponge (relaxation) layers: absorption in the momentum equation}

To prevent spurious reflections from the artificial domain edges, absorbing \emph{sponge} layers add a
Rayleigh-type linear damping to the horizontal momentum equation. For vertical layer $i$, the damping
force acts through the same vertical mass matrix $\VMij$ that appears in the acceleration term
\eqref{eq: horizontal momentum semi-discrete acceleration residual}, ensuring that all vertical modes
are damped consistently:
\begin{equation}
    \mathbf{Sp}_i^V \equiv \mu(\mathbf{x})\sumj \VMij\,\uj .
    \label{eq: sponge semi-discrete term}
\end{equation}
Weighting by $\vi$ and summing over $i$ as in the preceding subsections, and collecting the layers into the
multi-field vectors $\mathbf{U}$, $\mathbf{V}$ of
\eqref{eq: def U and V horizontal velocity multifield FE spaces}, the sponge contributes
\begin{equation}
    \boxed{\; R_{\text{sp}} = \sumi\left[\int_{\Oh}\mu(\mathbf{x})\left(\sumj\VMij\,\uj\right)\vi\;d\Oh\right] = \int_{\Oh}\mu(\mathbf{x})\left(\VmM\,\mathbf{U}\right)\cdot\mathbf{V}\;d\Oh \;}
    \label{eq: sponge residual contribution}
\end{equation}
The damping coefficient $\mu(\mathbf{x})\ge0$ is nonzero only inside the boundary layers and grows
\emph{quadratically} from zero at the inner edge of each layer to its maximum $\mu_{\max}$ at the
domain boundary. With up to four layers (left, right, bottom, top) of widths $w_L, w_R, w_B, w_T$ on a
rectangular domain $[x_0,x_1]\times[y_0,y_1]$, the profile is the pointwise maximum of the individual
edge ramps:
\begin{equation}
    \mu(\mathbf{x}) = \max\Big\{\, \mu_L, \mu_R, \mu_B, \mu_T \,\Big\},
    \label{eq: sponge profile max}
\end{equation}
with, for example, the left and bottom contributions
\begin{align}
    \mu_L(\mathbf{x}) &= \mu_{\max}\left(\frac{x_0 + w_L - x}{w_L}\right)^{2}
        \quad\text{for } x < x_0 + w_L, \qquad 0 \text{ otherwise}, \label{eq: sponge left}\\[1mm]
    \mu_B(\mathbf{x}) &= \mu_{\max}\left(\frac{y_0 + w_B - y}{w_B}\right)^{2}
        \quad\text{for } y < y_0 + w_B, \qquad 0 \text{ otherwise}, \label{eq: sponge bottom}
\end{align}
and symmetrically for the right ($x > x_1 - w_R$) and top ($y > y_1 - w_T$) edges. The argument of
each square is the normalised penetration depth into the layer, so $\mu$ ramps smoothly from $0$
(inner edge, no damping, no reflection) to $\mu_{\max}$ (outer wall, full damping). Taking the
\emph{maximum} rather than the sum in \eqref{eq: sponge profile max} clamps the overlapping corner
regions at $\mu_{\max}$ instead of $2\mu_{\max}$, avoiding an over-damped corner that would itself
reflect.

\paragraph{Design remarks.}
\begin{itemize}
    \item The sponge damps \emph{momentum} only; it does not appear in the continuity equation, so it
    removes kinetic energy while leaving mass conservation intact away from the generator.
    \item Because \eqref{eq: sponge residual contribution} is linear in $\mathbf{U}$, it adds directly
    and exactly to the momentum Jacobian block (no linearisation needed) and does not affect the
    effective mass matrix.
    \item The required strength scales with the wave regime: deep-water ($kd$ large) cases need larger
    $\mu_{\max}$ or wider layers than shallow-water cases, since a weak sponge lets amplitude build up
    over long runs.
\end{itemize}

\paragraph{Solver implementation.} The profile \eqref{eq: sponge profile max} is built by
\texttt{make\_sponge\_2D}$(\text{domain}, w_L, w_R, w_B, w_T, \mu_{\max})$, which returns the closure
$\mu(\mathbf{x})$. It is stored on the problem struct (\texttt{prob.mu\_sponge}) and, being
time-independent, is wrapped once into a static \texttt{CellField}; the term
\eqref{eq: sponge residual contribution} is appended to each momentum mode of the residual (and to the
Jacobian).

\subsubsection*{Augmented global residual}

Collecting \eqref{eq: wavemaker residual contribution} and
\eqref{eq: sponge residual contribution}, the complete residual passed to Gridap is the interior form
\eqref{eq: multifield general residual} together with the source and sponge contributions:
\begin{framed}
    \begin{multline}
        \text{Global Residual}_{\text{full}} = \text{Global Residual}\big|_{\eqref{eq: multifield general residual}}
        \;-\; \int_{\Oh} q\, S(\mathbf{x},t)\, d\Oh
        \;+\; \int_{\Oh} \mu(\mathbf{x})\left(\VmM\,\mathbf{U}\right)\cdot\mathbf{V}\; d\Oh .
        \label{eq: augmented global residual}
    \end{multline}
\end{framed}
The wavemaker term contributes only to the load vector (right-hand side), being independent of the
unknowns; the sponge term contributes symmetrically to both residual and Jacobian. Neither disturbs
the tensor-product separation of Section~7--8: the wavemaker is a purely horizontal forcing, and the
sponge factors as a horizontal scalar field $\mu(\mathbf{x})$ times the vertical mass matrix $\VmM$.




\section{Horizontal Discretisation and the Full Discrete Weak Form}
 
Having eliminated the vertical dimension $\sigma$ through the Galerkin weighted residual method applied along the water column in the previous sections, the system \eqref{eq: semi-discrete mass continuity residual}--\eqref{eq: semi-discrete horizontal momentum residual} now consists of $\Ns+2$ depth-integrated equations posed over the 2D horizontal domain $\Omega_h \subset \mathbb{R}^2$.
Each unknown is a purely 2D field: the $\Ns+1$ nodal horizontal velocity vectors $\uj(\mathbf{x},t)$ (for $j=0,\ldots,\Ns$) and the total water depth $H(\mathbf{x},t)$.
In this section we apply a second FE discretisation, now over $\Omega_h$, to reach the final fully discrete algebraic system.


 
\subsection{Horizontal Finite Element Approximation}
 
We introduce a partition of the horizontal domain $\Omega_h$ into $N_e$ non-overlapping elements.
Over this mesh we define a set of $\Nh+1$ scalar basis functions $\psi_n(\mathbf{x})$ of arbitrary polynomial order $q$, with $n=0,1,\ldots,\Nh$.
Consistent with the Galerkin methodology, the same basis is used for both trial and test functions.
 
Each nodal velocity field $\uj(\mathbf{x},t)$ and the total depth $H(\mathbf{x},t)$ are approximated as:
\begin{subequations}
    \begin{equation}
    \uj(\mathbf{x},t) = \begin{pmatrix} u_j(\mathbf{x}, t) \\ v_j(\mathbf{x}, t) \end{pmatrix} \equiv  \sumNhn \begin{pmatrix} \hat{u}_{jn}(t) \\ \hat{v}_{jn}(t) \end{pmatrix} \psi_n(\mathbf{x}) = \sumNhn \uhjn(t) \,\psi_n(\mathbf{x})
    \label{eq: horizontal velocity FE approximation}
\end{equation}
and
\begin{equation}
    H(\mathbf{x},t) \equiv \sumNhm \Hm(t) \,\psi_m(\mathbf{x})
    \label{eq: horizontal H FE approximation}
\end{equation}
\label{eq: horizontal FE approximations}
\end{subequations}
where $\uhjn = \begin{pmatrix} \hat{u}_{jn} & \hat{v}_{jn} \end{pmatrix} \in \mathbb{R}^2$ and $\Hm \in \mathbb{R}$ are the unknown time-dependent nodal degrees of freedom for the horizontal velocity at vertical node $j$ and horizontal node $n$, and for the total depth at horizontal node $m$, respectively.
 
We introduce horizontal test functions $\psi_s(\mathbf{x})$, $s=0,\ldots,\Nh$, drawn from the same space, and apply the Galerkin Weighted Residual Method over $\Omega_h$: multiply every equation of the system by $\psi_s(\mathbf{x})$ and integrate over the horizontal domain.
By the separation structure of the system, the resulting fully discrete equations take the form of tensor products between purely horizontal matrices (assembled over $\Omega_h$ in this section) and the purely vertical tensors $\VmM$, $\VtM$, $\VtG$, $\VmA$, $\VmK$, $\VtA$, $\VtK$ introduced in the previous sections.\\
 
Starting with the mass continuity equation \eqref{eq: semi-discrete mass continuity residual}, we multiply by $\psi_s(\mathbf{x})$ and integrate over $\Omega_h$:
$$\int_{\Omega_h} \left[ \ddt{H}\, \psi_s  + \sumj \Phi_j \nablah \cdot \left( H \uj\right) \psi_s \right] \, d\Omega = 0$$
We substitute the horizontal approximations \eqref{eq: horizontal FE approximations} into both terms and pull the time and vertical dependent terms out of the integrals since they do not depend on $\mathbf{x}$:
\begin{itemize}
    \item \textbf{Time derivative:}
    $$\int_{\Omega_h} \ddt{H}\,\psi_s\,d\Omega = \sumNhm \frac{d\Hm}{dt}\underbrace{\int_{\Omega_h}\psi_s\,\psi_m\,d\Omega}_{\HMsm}$$
    where we have isolated the horizontal mass matrix $\HmM$ with components:
    \begin{equation}
        \boxed{\HMsm \equiv \int_{\Omega_h}\psi_s\,\psi_m\,d\Omega} 
    \end{equation}
 
    \item \textbf{Divergence term:} Since $H$ and $\uj$ are both expanded over the same scalar basis, the bilinear product $H\uj$ in the divergence involves a convolution of two nodal coefficient series.
    To avoid this non-standard double sum in the divergence of a product, we integrate the divergence term by parts so that only first derivatives of the basis appear:
    \begin{equation}
        \int_{\Omega_h}\nablah\cdot(H\uj)\,\psi_s\,d\Omega = \oint_{\partial\Omega_h}(H\uj\cdot\mathbf{n})\,\psi_s\,ds - \int_{\Omega_h} H\uj\cdot\nablah\psi_s\,d\Omega
        \label{eq: continuity IBP}
    \end{equation}
    The boundary integral is handled by prescribing appropriate inflow/outflow conditions on $\partial\Omega_h$; on solid walls it vanishes since $\uj\cdot\mathbf{n}=0$.
    The interior integral is now a standard bilinear form.  Substituting \eqref{eq: horizontal FE approximations} yields:
    \begin{equation}
        \int_{\Omega_h} H\uj\cdot\nablah\psi_s\,d\Omega
        = \sumNhn\sumNhm \Hm \uhjn \cdot \underbrace{\int_{\Omega_h} \psi_m\,\psi_n\,\nablah\psi_s\,d\Omega}_{\HDsnma}
        \label{eq: continuity interior integral 1}
    \end{equation}
    where $\HttD$ is a 4th-order static horizontal tensor, for which the components are given by:
    \begin{equation}
        \boxed{\mathbb{D}^H_{snm, a} \equiv (\HDsnma)_a \equiv \int_{\Omega_h} \psi_m\,\psi_n\,\big(\nablah\psi_s \big)_a\,d\Omega \equiv \int_{\Omega_h} \psi_m \psi_n \ddxa{\psi_s} \, d\Omega}
        \label{eq: def horizontal mass continuity divergence operator}
    \end{equation}
    In the following, such to alleviate the notation and compress the degrees of freedom in each horizontal node, we consider the contraction $\uhjn \cdot \HDsnma$ and drop the horizontal component index without loss of generality. 
\end{itemize}

Collecting these terms, the discrete mass continuity equation at horizontal test node $s$ is:
\begin{equation}
\boxed{\sumNhm \frac{d\Hm}{dt}\,\HMsm - \sumj\Phi_j \sumNhn\sumNhm \Hm\uhjn \cdot \HDsnma = -\sumj\Phi_j\,\mathbf{F}_{sj}}
    \label{eq: fully discrete mass continuity}
\end{equation}
where $\mathbf{F}$ is the matrix collecting the prescribed boundary fluxes at layer $j$ and boundary horizontal node $s$ with components
\begin{equation}
    \mathbf{F}_{sj} \equiv \oint_{\partial\Omega_h}(H\uj\cdot\mathbf{n})\,\psi_s\,ds  
    \label{eq: horizontal mass continuity boudary fluxes}
\end{equation}

We take this instant to complete the notation and introduce the following arrays of unknowns:
\begin{equation}
    \mathbf{H} \equiv \begin{bmatrix} 
    \hat{H}_0 \\[1mm]
    \hat{H}_1 \\[1mm] 
    \vdots \\[1mm]
    \hat{H}_{\Nh -1}\\[1mm]
    \hat{H}_{\Nh}
    \end{bmatrix} \qquad \text{and} \qquad  
    \mathbf{U}_j \equiv \begin{bmatrix} 
    \big(\hat{\mathbf{u}}_{j,0}\big)^T \\[1mm]
    \big(\hat{\mathbf{u}}_{j,1}\big)^T \\[1mm]
    \vdots \\[1mm]
    \big(\hat{\mathbf{u}}_{j,\Nh -1}\big)^T\\[1mm]
    \big(\hat{\mathbf{u}}_{j,\Nh}\big)^T
    \end{bmatrix} = \begin{bmatrix}
    \hat{u}_{j,0} & \hat{v}_{j,0} \\[1.5mm]
    \hat{u}_{j,1} & \hat{v}_{j,1} \\[1.5mm]
    \vdots & \vdots \\[1.5mm]
    \hat{u}_{j,\Nh -1} & \hat{v}_{j,\Nh -1} \\[1.5mm]
    \hat{u}_{j,\Nh} & \hat{v}_{j,\Nh} 
    \end{bmatrix} = \begin{bmatrix}
        \hat{\mathbf{u}}_j & \hat{\mathbf{v}}_j
    \end{bmatrix}
    \label{eq: def array of horizontal unknowns}
\end{equation}
where $\mathbf{H}$ corresponds to the vector of horizontal nodal water depths, and $\mathbf{U}_j$ corresponds to the matrix containing the horizontal nodal velocity degrees of freedom at layer $j$.

Hence, using this compact notation we are able to write horizontal mass continuity equation as follows:
$$\HmM \dot{\Hvec} - \sumj\Phi_j \, \Hvec\cdot(\HttD :\Uj) = -\mathbf{F}\,\boldsymbol{\Phi}$$
where the sums are translated into contractions of indices and $\mathbf{F}$ is the $(\Nh+1)\times(\Ns+1)$ matrix of boundary fluxes $\mathbf{F}_{sj}$.

An alternative process can be established in order to avoid excessive contractions.
This is done by defining depth-weighted horizontal operator matrices, which means including the total depth $H(\mathbf{x},t)$ when computing the operator matrix terms, and calculated on-the-fly at each time step using the current state of $\mathbf{H}$. 

For instance, retaking the interior integral we can group the terms in this case:
\begin{equation}
    \int_{\Omega_h} H\uj\cdot\nablah\psi_s\,d\Omega =  \int_{\Omega_h} H\left(\sumNhn\uhjn \psi_n\right) \cdot \nablah\psi_s\,d\Omega = \sumNhn \uhjn \cdot \underbrace{\int_{\Omega_h} H(\mathbf{x},t)\,\psi_n\,\nablah\psi_s\,d\Omega}_{\HDdwsn(H)} 
    \label{eq: continuity interior integral 2}
\end{equation}
where the depth-weighted horizontal divergence operator 3rd-order tensor $\HtDdw(H)$ has components for a horizontal test node $s$ and trial node $n$ given by expression:
\begin{equation}
    \boxed{\HtDdw_{sn,a}(H) \equiv \big(\HDdwsn(H)\big)_a \equiv \int_{\Omega_h} H(\mathbf{x}, t) \, \psi_n \big( \nabla_h \psi_s \big)_a\, d\Omega \equiv \int_{\Omega_h} H(\mathbf{x}, t) \, \psi_n \ddxa{\psi_s}\, d\Omega}
    \label{eq: def depth weighted horizontal mass continuity divergence operator}
\end{equation}
Hence, observing expressions \eqref{eq: continuity interior integral 1} and \eqref{eq: continuity interior integral 2} we can naturally identify :
$$\int_{\Omega_h} H\uj\cdot\nablah\psi_s\,d\Omega
        = \sumNhn\sumNhm \Hm \uhjn \cdot \HDsnma = \sumNhn \uhjn \cdot \HDdwsn(H)$$
\begin{equation}
    \implies \quad \boxed{\HtDdwH = \sumNhm \Hm\, \HDsnma =  \Hvec^T \cdot \HttD = \HttD \cdot \Hvec}
    \label{eq: def depth weighted horizontal mass continuity divergence operator functional}
\end{equation}
Using definition \eqref{eq: def depth weighted horizontal mass continuity divergence operator functional}, the horizontal mass continuity equation can be written now as
$$\HmM \dot{\Hvec} - \sumj\Phi_j \, (\HtDdwH :\Uj) = -\mathbf{F}\,\boldsymbol{\Phi}$$

Now, we bring back the vertical layer coupling. 
As far as the mass continuity equation is concerned, the coupling interaction is represented by the product with factors $\Phi_j$ for each layer $j$. 

We proceed by stacking all vertical layer unknowns into a single global velocity block matrix $\Ublock$ with size $(\Ns+1)(\Nh+1 )\times 2$ the following general form:
\begin{equation}
    \Ublock \equiv \begin{bmatrix}
    \mathbf{U}_0 \\[1mm]
    \mathbf{U}_1 \\[1mm]
    \vdots \\[1mm]
    \mathbf{U}_{\Ns}
    \end{bmatrix}
    \label{eq: def vertically stacked horizontal velocities}
\end{equation}
where each entry $(\Ublock)_j = \Uj$ corresponds to the horizontal velocity degrees of freedom for layer $j=0,...,\Ns$.  

With this structure, the horizontal mass continuity equation vertical layer dependencies can be fully contracted by taking advantage of the Kronecker product ($\otimes$).
This reduces the entire mass continuity equation across all vertical layers and horizontal nodes down to a single algebraic system of $\Nh+1$ equations:   
\begin{equation}
    \boxed{\HmM \dot{\Hvec} - \big(\boldsymbol{\Phi}^T \otimes \HtDdwH \big): \Ublock =  -\mathbf{F} \, \boldsymbol{\Phi}}
    \label{eq: fully compresses mass continuity equation}
\end{equation}
where we have introduced (as in Section 8) the vector of stacked depth-integrated basis profiles
\begin{equation}
    \boldsymbol{\Phi} \equiv \begin{bmatrix}
    \Phi_0 & \Phi_1 & \dots & \Phi_{\Ns}
    \end{bmatrix}^T
    \label{eq: def array Phi discrete}
\end{equation}




 
\subsection{Discretisation of the Horizontal Momentum Equation}

Regarding the horizontal momentum residual \eqref{eq: semi-discrete horizontal momentum residual}, we discretise it now introducing \eqref{eq: horizontal FE approximations} and applying the Weighted Residual Method at each vertical layer $i$: since the momentum residual is a 2D \emph{vector} equation, weighting by the scalar test function $\psi_s(\mathbf{x})$ and integrating over $\Omega_h$ yields \emph{two} scalar equations per horizontal node (one per horizontal component). We keep the vector notation throughout and make the component index explicit only where a basis derivative forces it:
\begin{multline}
    \int_{\Omega_h}\left[\sumj H\ddt{\uj}\VMij
    + \sumk\sumj\Big(H\VMikj(\uk\cdot\nablah\uj) + (\nablah\cdot[H\uk])\uj\,\VGikj\Big)
    + gH\nablah\eta\,\Phi_i\right]\psi_s\,d\Omega \\
    = \int_{\Omega_h}\Bigg[ -\nablah\left( H^2 \sumj \mathcal{L}_j\cdot\VPij + H^2 \sumk\sumj \mathcal{N}_{kj}\cdot\VPikj \right) \\
    + H\sumj\Big(\nablah h\,(\mathcal{L}_j\cdot\VAij)
      + \nablah H\,(\mathcal{L}_j\cdot\VKij)\Big) \\
    + H\sumk\sumj\Big(\nablah h\,(\mathcal{N}_{kj}\cdot\VAikj)
      + \nablah H\,(\mathcal{N}_{kj}\cdot\VKikj)\Big)\Bigg]\psi_s\,d\Omega
    \label{eq: horizontal WRM momentum}
\end{multline}
Because $H$, $\uj$, $\uk$, $\nablah h$, $\nablah H$ all live in the horizontal domain only, and because the vertical tensors $\VMij$, $\VMikj$, etc., are real-valued scalars (or component vectors) that pass outside the $\Omega_h$ integral, we treat each term in sequence below. For each term we follow the same two-step pattern used for the mass continuity equation: first the fully \emph{static} form, isolating a state-independent horizontal tensor assembled once; then the \emph{state-weighted} alternative, a lower-order operator recomputed on-the-fly from the current state of $\Hvec$ (and, for the advection, of $\Uj$), which is the form a practical solver assembles.

\subsubsection*{Acceleration term}

In the acceleration term we substitute \eqref{eq: horizontal FE approximations} and pull outside of the integral all factors which do not vary horizontally, yielding:
\begin{equation}
    \mathbf{Acc}_{i,s}^H = \int_{\Omega_h} \left[\sumj H\ddt{\uj}\VMij\right]\psi_s\,d\Omega
    = \sumj \VMij \sumNhn\sumNhm \Hm\frac{d\uhjn}{dt} \underbrace{\int_{\Omega_h}\psi_s \,\psi_n\,\psi_m\,d\Omega}_{\HMsnm}
    \label{eq: discrete acceleration}
\end{equation}
where we isolated the fully symmetric 3rd-order horizontal mass tensor $\HtM$ with components:
\begin{equation}
    \boxed{\HMsnm \equiv \int_{\Omega_h}\psi_s\,\psi_n\,\psi_m\,d\Omega}
    \label{eq: def horizontal mass 3rd order}
\end{equation}

Similarly to what was done with $\HttD$ and $\HtDdwH$ in the previous subsection, we can lower the order of the operator involved by depth-weighting, which means including the total depth $H(\mathbf{x},t)$ when computing the operator, using the current state of $\mathbf{H}$.
Hence, we define the depth-weighted horizontal mass matrix $\HmMdwH$, recomputed at each time step:
\begin{equation}
    \boxed{\bar{\mathbf{M}}_{sn}^H(H) \equiv \int_{\Omega_h} H(\mathbf{x}, t) \, \psi_s \, \psi_n \, d\Omega = \sumNhm \Hm \, \HMsnm = \HtM \cdot \Hvec}
    \label{eq: def depth weighted horizontal mass}
\end{equation}
Using this notation, the horizontal summations collapse into a standard matrix-vector product per layer pair $(i,j)$:
$$\mathbf{Acc}_{i}^H = \sumj \VMij \, \HmMdwH \, \dot{\mathbf{U}}_j$$
Now, using the global velocity block matrix $\Ublock$ introduced in \eqref{eq: def vertically stacked horizontal velocities}, the complete acceleration term for the global system can be written elegantly without a single summation symbol as a Kronecker product:
\begin{equation}
    \mathbf{Acc}^H = \left( \VmM \otimes \HmMdwH \right) \dot{\Ublock}
    \label{eq: discrete acceleration Kronecker}
\end{equation}



\subsubsection*{Horizontal advection term}

The same process is carried out on the advection terms at layer $i$. We start with the horizontal advection contribution. Substituting \eqref{eq: horizontal FE approximations}, the advected gradient expands as $(\uk\cdot\nablah)\uj = \sumNhn\sumNhr \psi_n \big(\uhkn\cdot\nablah\psi_r\big)\,\uhjr$, so that:
\begin{align*}
    \mathbf{HAdv}_{i,s}^H &= \int_{\Omega_h} \left[\sumk\sumj H\VMikj(\uk\cdot\nablah\uj)\right]\psi_s\,d\Omega \\
    &= \sumk\sumj \VMikj \sumNhm \sumNhn \sumNhr \Hm \left(\uhkn \cdot \int_{\Omega_h} \psi_s\, \psi_m\, \psi_n\, \nablah \psi_r \,d\Omega \right) \uhjr \\
    &= \sumk\sumj \VMikj \sumNhm \sumNhn \sumNhr \Hm \left(\uhkn \cdot \mathbb{T}^H_{smnr}\right) \uhjr
\end{align*}
where the quadruple product of basis functions (one carrying the gradient) produces the static 4th-order, vector-valued horizontal advection tensor:
\begin{equation}
    \boxed{\mathbb{T}^H_{smnr} \equiv \int_{\Omega_h}\psi_s\,\psi_m\,\psi_n\,\nablah\psi_r\,d\Omega
    \qquad\Longleftrightarrow\qquad
    \big(\mathbb{T}^H_{smnr}\big)_a = \int_{\Omega_h}\psi_s\,\psi_m\,\psi_n\,\ddxa{\psi_r}\,d\Omega}
    \label{eq: def horizontal 4th order tensor}
\end{equation}

As before, the practical alternative is to \emph{state-weight} the operator: freezing the current fields $H$ and $\uk$ inside the integral collapses the four nested sums into one sparse matrix per vertical layer $k$, recomputed at each nonlinear iteration:
\begin{equation}
    \boxed{\bar{\mathbf{C}}^H_{sr}(H,\uk) \equiv \int_{\Omega_h} \psi_s \, H \left(\uk\cdot\nablah\psi_r\right) d\Omega}
    \qquad\implies\qquad
    \mathbf{HAdv}_{i}^H = \sumk\sumj \VMikj \, \bar{\mathbf{C}}^H(H,\uk)\, \mathbf{U}_j
    \label{eq: def state weighted advection operator}
\end{equation}

\subsubsection*{Vertical advection term}

For the transformed vertical advection contribution, the divergence factor expands with the product rule $\nablah\cdot(H\uk) = H\nablah\cdot\uk + \uk\cdot\nablah H$, i.e.
$$\nablah\cdot(H\uk) = \sumNhm\sumNhn \Hm \, \uhkn\cdot\big( \psi_m\nablah\psi_n + \psi_n\nablah\psi_m \big)$$
so that, together with $\uj = \sumNhr \uhjr\psi_r$:
\begin{align*}
    \mathbf{VAdv}_{i,s}^H &= \int_{\Omega_h}\left[\sumk\sumj(\nablah\cdot[H\uk])\,\uj\,\VGikj\right]\psi_s\,d\Omega \\
    &= \sumk\sumj \VGikj \sumNhm\sumNhn\sumNhr \Hm \Big[\uhkn\cdot\big(\mathbb{T}^H_{smrn} + \mathbb{T}^H_{snrm}\big)\Big] \uhjr
\end{align*}
where \emph{no new tensor is needed}: both pieces are index permutations of the same 4th-order tensor $\mathbb{T}^H$ defined in \eqref{eq: def horizontal 4th order tensor} (the gradient falling once on the velocity basis $\psi_n$ and once on the depth basis $\psi_m$).

The state-weighted alternative freezes the whole scalar field $\nablah\cdot(H\uk)$ from the current state, again one sparse matrix per layer $k$ per nonlinear iteration:
\begin{equation}
    \boxed{\bar{\mathbf{S}}^H_{sr}(H,\uk) \equiv \int_{\Omega_h} \psi_s \, \big(\nablah\cdot[H\uk]\big)\,\psi_r \, d\Omega}
    \qquad\implies\qquad
    \mathbf{VAdv}_{i}^H = \sumk\sumj \VGikj \, \bar{\mathbf{S}}^H(H,\uk)\, \mathbf{U}_j
    \label{eq: def state weighted vertical advection operator}
\end{equation}
The pair $\big(\bar{\mathbf{C}}^H(H,\uk), \bar{\mathbf{S}}^H(H,\uk)\big)$ contracted with the small vertical tensors $\VtM,\VtG$ is precisely the matrix-free ``V$\otimes$H'' route implemented in the production solver: the enormous static tensor $\mathbb{T}^H$ is never stored.
 
\subsubsection*{Gravity term}

For the gravity term, we use $\nablah\eta = \nablah H - \nablah h$ (since $h$ is prescribed). Both pieces are directly expandable — $\nablah H = \sumNh{\gamma}\hat{H}_\gamma\nablah\psi_\gamma$ contains only first derivatives of the basis, so \emph{no} integration by parts is required:
$$\mathbf{Grav}_{i,s}^H = g\Phi_i\int_{\Omega_h}H\big(\nablah H - \nablah h\big)\,\psi_s\,d\Omega
= g\Phi_i\left[ \sumNhm\sumNh{\gamma}\Hm\hat{H}_\gamma\,\mathbb{E}^H_{sm\gamma} - \sumNhm\Hm\,\mathbb{B}^h_{sm} \right]$$
with the static 3rd-order gradient tensor and the (bed-slope--weighted) 2nd-order bathymetry tensor:
\begin{equation}
    \boxed{\mathbb{E}^H_{sm\gamma} \equiv \int_{\Omega_h}\psi_s\,\psi_m\,\nablah\psi_\gamma\,d\Omega}
    \qquad\qquad
    \boxed{\mathbb{B}^h_{sm} \equiv \int_{\Omega_h}\psi_s\,\psi_m\,\nablah h\,d\Omega}
    \label{eq: def horizontal gravity tensors}
\end{equation}
Note that $\mathbb{E}^H$ is the same object as the continuity tensor $\HttD$ with permuted indices, $\mathbb{E}^H_{sm\gamma} = \HDsnma\big|_{s\to\gamma,\,n\to m,\,m\to s}$: the derivative simply falls on a trial function instead of the test function.

\medskip
\noindent\textit{Remark (practical form).} An equivalent \emph{energy} form uses $gH\nablah H = \tfrac{g}{2}\nablah(H^2)$ and integrates by parts onto the test function. In a discrete setting where the resulting boundary flux is dropped (natural boundary conditions at open boundaries), the constant hydrostatic baseline must then be subtracted, $\nablah(H^2)\to\nablah(H^2-d^2)$ with $d$ the still-water depth, so that the discrete rest state $\eta\equiv 0$ carries exactly zero net force; otherwise a spurious wall force accelerates the fluid from rest. The direct form above does not suffer from this issue but places the derivative on the (quadratic) trial expansion.
 
\subsubsection*{Right-hand side: leading pressure term (dispersion)}

The retained leading gradient contributes $-\int_{\Omega_h}\nablah\big(H^2\,\mathcal{S}_i\big)\,\psi_s\,d\Omega$ per layer $i$, where we abbreviate the scalar pressure moment
$$\mathcal{S}_i \equiv \sumj \mathcal{L}_j\cdot\VPij + \sumk\sumj \mathcal{N}_{kj}\cdot\VPikj \;.$$
Integrating by parts onto the test function (the boundary integral vanishes on essential/solid-wall boundaries or behind sponge layers):
$$-\int_{\Omega_h}\nablah\big(H^2\mathcal{S}_i\big)\,\psi_s\,d\Omega = \int_{\Omega_h} H^2\,\mathcal{S}_i\,\nablah\psi_s\,d\Omega - \oint_{\partial\Omega_h} H^2\,\mathcal{S}_i\,\psi_s\,\mathbf{n}\,ds$$
The dominant, dispersion-carrying piece is the third component of the linear part, $(\mathcal{L}_j)_3(\VPij)_3 = -\nablah\cdot(H\partial_t\uj)\,(\VPij)_3$ with $(\VPij)_3 = -B_{ij}$. Expanding $\nablah\cdot(H\partial_t\uj)$ with the product rule as in the vertical advection term:
$$\int_{\Omega_h} H^2\big[\nablah\cdot(H\ddt{\uj})\big]\nablah\psi_s\,d\Omega
= \sumNhm\sumNh{\gamma}\sumNh{\delta}\sumNhn \Hm\hat{H}_\gamma\hat{H}_\delta\,\Big[\frac{d\uhjn}{dt}\cdot\big(\mathbb{Q}^H_{sm\gamma\delta n} + \mathbb{Q}^H_{sm\gamma n\delta}\big)\Big]$$
with the 5th-order static family $\mathbb{Q}^H_{sm\gamma\delta n} \equiv \int_{\Omega_h}\nablah\psi_s\otimes\nablah\psi_n\,\psi_m\psi_\gamma\psi_\delta\,d\Omega$ (a 2$\times$2-matrix-valued integral). Denoting by $\mathbf{P}_{i}^H$ this whole leading-pressure contribution \emph{as it stands on the right-hand side} of the momentum equation (after the IBP above, flux dropped), in the flat-bed linearised limit ($H\to d$ constant) it collapses onto the classical stiffness tensor:
\begin{equation}
    \boxed{\mathbb{K}^H_{sn} \equiv \int_{\Omega_h}\nablah\psi_s\otimes\nablah\psi_n\,d\Omega}
    \qquad\implies\qquad
    \mathbf{P}_{i}^H\Big|_{\text{flat,lin}} = d^3\sumj B_{ij}\,\mathbb{K}^H\,\dot{\mathbf{U}}_j
    \label{eq: def horizontal stiffness tensor}
\end{equation}
which is the discrete frequency-dispersion operator of the model (recall $B_{ij} = -\int_0^1\vphij\varphi_i\,d\sigma \le 0$, so the term opposes the acceleration, as a dispersive pressure must). In practice one never assembles $\mathbb{Q}^H$: the state-weighted operator
\begin{equation}
    \boxed{\bar{\mathbf{P}}^H_{sn}(H) \equiv \int_{\Omega_h} H^2\,\nablah\psi_s\otimes\big( H \nablah\psi_n + \psi_n \nablah H \big)\, d\Omega}
    \label{eq: def state weighted leading pressure operator}
\end{equation}
is recomputed from the current $\Hvec$ and gives
$$\mathbf{P}_i^H = \sumj B_{ij}\,\bar{\mathbf{P}}^H(H)\,\dot{\mathbf{U}}_j + \big(\text{remaining } \VPij, \VPikj \text{ components, same pattern}\big).$$

\subsubsection*{Right-hand side: linear slope pressure terms}

The linear slope-pressure RHS at vertical layer $i$ is:
$$\mathbf{LRHS}_{i,s} = \int_{\Omega_h}H\sumj\Big[\nablah h\,(\mathcal{L}_j\cdot\VAij) + \nablah H\,(\mathcal{L}_j\cdot\VKij)\Big]\psi_s\,d\Omega$$
Recalling that $\mathcal{L}_j = \big(-\partial_t\uj\cdot\nablah h,\; \partial_t\uj\cdot\nablah H,\; -\nablah\cdot(H\partial_t\uj)\big)^T$, every component is \emph{linear} in $\dot{\uj} = \partial_t\uj$: in nodal form,
$$(\mathcal{L}_j)_1 = -\sumNhn \big(\tfrac{d\uhjn}{dt}\cdot\nablah h\big)\psi_n, \qquad
(\mathcal{L}_j)_2 = \sumNhn \big(\tfrac{d\uhjn}{dt}\cdot\nablah H\big)\psi_n,$$
$$(\mathcal{L}_j)_3 = -\sumNhn \frac{d\uhjn}{dt}\cdot\big( H\nablah\psi_n + \psi_n\nablah H \big)$$
where the third component used the product rule $\nablah\cdot(H\dot{\uj}) = H\nablah\cdot\dot{\uj} + \dot{\uj}\cdot\nablah H$ — only first derivatives of the basis appear, so \emph{no} integration by parts is needed anywhere in this block.

Substituting these expansions produces, per component $\ell$, integrals of the pattern $\int\psi_s\,H\,(\text{slope})\otimes(\text{coefficient of }\tfrac{d\uhjn}{dt})\,d\Omega$. Following the state-weighted approach, we collect them into six sparse, $2\times2$-block-valued matrices, recomputed from the current state ($\nablah H$ evaluated from $\Hvec$, $\nablah h$ prescribed):
\begin{subequations}
\begin{align}
    \boxed{\bar{\mathbb{A}}^{(1)}_{sn} \equiv -\int_{\Omega_h}\psi_s\,\psi_n\,H\;\nablah h\otimes\nablah h\;d\Omega}\\[2pt]
    \boxed{\bar{\mathbb{A}}^{(2)}_{sn} \equiv \int_{\Omega_h}\psi_s\,\psi_n\,H\;\nablah h\otimes\nablah H\;d\Omega}\\[2pt]
    \boxed{\bar{\mathbb{A}}^{(3)}_{sn} \equiv -\int_{\Omega_h}\psi_s\,H\;\nablah h\otimes\big(H\nablah\psi_n + \psi_n\nablah H\big)\,d\Omega}
\end{align}
\label{eq: def state weighted linear pressure operators}
\end{subequations}
and $\bar{\mathbb{K}}^{(\ell)}_{sn}$ identically with the prefactor $\nablah h$ replaced by $\nablah H$. The full linear slope-pressure RHS is then the compact double contraction:
\begin{equation}
    \mathbf{LRHS}_{i} = \sumj\sum_{\ell=1}^{3}\Big[(\VAij)_\ell\,\bar{\mathbb{A}}^{(\ell)} + (\VKij)_\ell\,\bar{\mathbb{K}}^{(\ell)}\Big]\dot{\mathbf{U}}_j
    \label{eq: discrete linear RHS}
\end{equation}
Fully static expansions exist (substituting $H=\sum_m\Hm\psi_m$ and $\nablah H = \sum_\gamma\hat H_\gamma\nablah\psi_\gamma$ everywhere generates 3rd- to 5th-order tensors of the $\mathbb{T}^H$/$\mathbb{Q}^H$ families) but proliferate high-order objects with no practical benefit over the state-weighted form.
 
\subsubsection*{Right-hand side: non-linear pressure terms}

The non-linear RHS involves $\mathcal{N}_{kj}$ contracted with $\VAikj$ and $\VKikj$ (and $\VPikj$ in the leading term). Since $\mathcal{N}_{kj}$ is \emph{quadratic} in $(\uk,\uj)$ and carries factors of $H$, substituting the horizontal expansion \eqref{eq: horizontal FE approximations} generates cubic combinations of nodal coefficients: fully static expansions would require 5th- and 6th-order tensors and are not worth storing. The natural discrete treatment is again state-weighted, per component $c = 1,\dots,8$ and per layer pair $(k,j)$, following exactly the advection pattern: evaluate the scalar building blocks $\nablah\cdot(H\uk)$, $\uk\cdot\nablah h$, $\uk\cdot\nablah H$ from the current state, and integrate them against $\psi_s$ with the $\nablah h$/$\nablah H$ prefactors.

Two structural observations, however, must be respected:
\begin{itemize}
    \item Components $c = 6,7,8$ of $\mathcal{N}_{kj}$ contain only \emph{first} derivatives of the unknowns (products of the scalar building blocks above, divided by $H$ — the $1/H$ cancels against the overall $H$ prefactor of the RHS). They discretise directly, with no integration by parts.
    \item Components $c = 1,\dots,5$ contain \emph{second} derivatives of the unknowns ($\nablah(\nablah\cdot[H\uk])$ and, through $\nablah(\uj\cdot\nablah H)$, the Hessian of the free surface). On a $C^0$ FE basis these are not directly representable. For the $\nablah h$-prefactored half, one integration by parts onto the test function removes the offending derivative exactly (the only second derivative left is the \emph{prescribed} bed Hessian $\nablah^2 h$). For the $\nablah H$-prefactored half the second derivative of the unknown is irreducible: it requires either an $L^2$ projection of the inner divergence onto the FE space before differentiating, or an auxiliary mixed field. These non-linear pressure contributions are $O(A^3)$ in wave amplitude and are only needed for finite-amplitude harmonic generation or strong shoaling.
\end{itemize}
Symbolically, the discrete non-linear RHS at layer $i$ reads
\begin{equation}
    \mathbf{NLRHS}_{i} = \sumk\sumj\sum_{c=1}^{8}\Big[(\VAikj)_c\,\bar{\mathbb{A}}^{(c)}_{\mathcal{N}}(H,\uk,\uj) + (\VKikj)_c\,\bar{\mathbb{K}}^{(c)}_{\mathcal{N}}(H,\uk,\uj)\Big]
    \label{eq: discrete nonlinear RHS}
\end{equation}
where $\bar{\mathbb{A}}^{(c)}_{\mathcal{N}},\bar{\mathbb{K}}^{(c)}_{\mathcal{N}}$ are the state-weighted vectors $\int_{\Omega_h}\psi_s\,H\,(\nablah h\ \text{resp.}\ \nablah H)\,(\mathcal{N}_{kj})_c\,d\Omega$ assembled from the current state with the componentwise treatment described above.
 
\subsection{Final Fully Discrete Weak Form}

Assembling all contributions from the two previous subsections, the final fully discrete system takes the form of contractions between the small dense \emph{vertical} tensors of Sections~6--7 and the sparse \emph{horizontal} operators of this section (static where practical, state-weighted otherwise).
It consists of $(2\Ns+3)(\Nh+1)$ scalar algebraic equations — one continuity equation plus two momentum components per horizontal node and vertical layer — in the unknown nodal coefficients $\{\uhjn\}$ and $\{\Hm\}$. Written per horizontal test node $s$ (momentum in 2D vector form) using the state-weighted operators:

\begin{framed}
\noindent\textbf{Discrete mass continuity} at horizontal test node $s$:
\begin{equation}
    \sumNhm\frac{d\Hm}{dt}\,\HMsm
    - \sumj\Phi_j\,\sumNhn \uhjn\cdot\HDdwsn(H)
    = -\sumj\Phi_j\,\mathbf{F}_{sj}
    \label{eq: final discrete mass continuity}
\end{equation}

\bigskip
\noindent\textbf{Discrete horizontal momentum} at horizontal test node $s$, vertical layer $i$:
\begin{multline}
    \underbrace{\sumj\VMij\,\bar{\mathbf{M}}^H(H)\,\dot{\mathbf{U}}_j}_{\text{Acceleration}}
    +\underbrace{\sumk\sumj\Big[\VMikj\,\bar{\mathbf{C}}^H(H,\uk) + \VGikj\,\bar{\mathbf{S}}^H(H,\uk)\Big]\mathbf{U}_j}_{\text{Advection}}\\
    +\underbrace{g\Phi_i\Big[\sumNhm\sumNh{\gamma}\Hm\hat{H}_\gamma\,\mathbb{E}^H_{sm\gamma} - \sumNhm\Hm\,\mathbb{B}^h_{sm}\Big]}_{\text{Gravity}} \\[6pt]
    = \underbrace{\sumj \Big[ (\VPij)_3\,(-1)\,\bar{\mathbf{P}}^H(H)\,\dot{\mathbf{U}}_j + \dots\Big]}_{\text{Leading pressure (dispersion), \eqref{eq: def state weighted leading pressure operator}}}
    + \underbrace{\sumj\sum_{\ell=1}^{3}\Big[(\VAij)_\ell\,\bar{\mathbb{A}}^{(\ell)} + (\VKij)_\ell\,\bar{\mathbb{K}}^{(\ell)}\Big]\dot{\mathbf{U}}_j}_{\text{Linear slope pressure, \eqref{eq: discrete linear RHS}}}\\
    + \underbrace{\mathbf{NLRHS}_{i}}_{\text{Non-linear pressure, \eqref{eq: discrete nonlinear RHS}}}
    \;+\; \mathcal{F}_{s,i}^{\text{mom}}
    \label{eq: final discrete momentum}
\end{multline}
\end{framed}

\noindent where $\mathcal{F}_{s,i}^{\text{mom}}$ collects all boundary flux contributions arising from the integration-by-parts steps applied throughout (they vanish on solid walls and behind sponge layers), and the leading-pressure bracket abbreviates the $(\VPij)_3$-dominant form of \eqref{eq: def state weighted leading pressure operator} — the sign $(-1)$ makes the $B_{ij} = -(\VPij)_3$ convention explicit — plus the remaining $\VPij$ and $\VPikj$ components.
 
\subsection{Structure of the System and Tensor Product Interpretation}
 
The system \eqref{eq: final discrete mass continuity}--\eqref{eq: final discrete momentum} has a clear two-level tensor product structure that reflects the original two-step discretisation:
 
\begin{enumerate}
    \item \textbf{Vertical tensors} $\VMij$, $\VMikj$, $\VGikj$, $\VAij$, $\VKij$, $\VPij$, $\VAikj$, $\VKikj$, $\VPikj$, $\Phi_j$ are computed once from the vertical FE mesh $\{\sigma_j\}$ and stored. They are dense but small: their size scales as $O(\Ns^2)$ or $O(\Ns^3)$.

    \item \textbf{Horizontal operators} come in two families. The purely \emph{static} tensors $\HmM$, $\HtM$, $\HttD$, $\mathbb{E}^H$, $\mathbb{B}^h$, $\mathbb{K}^H$ (and, formally, $\mathbb{T}^H$, $\mathbb{Q}^H$) are assembled from the horizontal mesh once and are large but sparse: their non-zero entries scale as $O(\Nh)$ (each basis function interacts only with its neighbours). The \emph{state-weighted} operators $\HmMdwH$, $\HtDdwH$, $\bar{\mathbf{C}}^H(H,\uk)$, $\bar{\mathbf{S}}^H(H,\uk)$, $\bar{\mathbf{P}}^H(H)$, $\bar{\mathbb{A}}^{(\ell)}$, $\bar{\mathbb{K}}^{(\ell)}$ are equally sparse matrices re-assembled from the current state at each time step / nonlinear iteration; they trade a small assembly cost for a drastic reduction in tensor order.

    \item The fully discrete residual at degree of freedom $(s,i)$ is always a \textbf{contraction of one small vertical tensor entry with one sparse horizontal operator entry}:
    \begin{equation*}
        \text{Residual}_{s,i} = \sum_{j,k,\ldots} (\text{vertical tensor})_{i,k,j}\;\times\;(\text{horizontal operator})_{s,n,\ldots}\;\times\;(\text{nodal DOFs})
    \end{equation*}
    This structure allows the assembly of the global system to be performed as a loop over vertical layer pairs $(i,j)$ and triplets $(i,k,j)$, reusing the horizontal operators of the current step. It is also exactly the structure that a high-level FE toolbox (Section 8) exploits automatically when handed the continuous weak form.
\end{enumerate}
 
The complete list of pre-computable static tensors and their roles is summarised in Table~\ref{tab: all tensors}.
 
\begin{table}[H]
\centering
\renewcommand{\arraystretch}{1.45}
\begin{tabular}{llll}
\toprule
Tensor & Definition & Order & Role \\
\midrule
\multicolumn{4}{l}{\textit{Vertical tensors (from Sections 3--7)}}\\
$\VMij$ & $\int_0^1\phi_i\phi_j\,d\sigma$ & $2$ & Vertical mass \\
$\Phi_j$ & $\int_0^1\phi_j\,d\sigma$ & $1$ & Depth-average weight \\
$\VMikj$ & $\int_0^1\phi_i\phi_k\phi_j\,d\sigma$ & $3$ & Vertical advection mass \\
$\VGikj$ & $\int_0^1(\sigma\Phi_k-\varphi_k)\phi'_j\phi_i\,d\sigma$ & $3$ & Transformed vert.\ advection \\
$\VAij$ & $\int_0^1\phi_i\,\theta_j\,d\sigma$ & $2$ & Linear pressure (bed slope) \\
$\VKij$ & $\int_0^1(\Pi_j-\sigma\theta_j)\phi_i\,d\sigma$ & $2$ & Linear pressure (surface slope) \\
$\VAikj$ & $\int_0^1\Theta_{kj}\phi_i\,d\sigma$ & $3$ & Non-linear pressure \\
$\VKikj$ & $\int_0^1(\Lambda_{kj}-\sigma\Theta_{kj})\phi_i\,d\sigma$ & $3$ & Non-linear pressure \\
$\VPij$ & $\int_0^1\theta_j\,\varphi_i\,d\sigma$ & $2$ & Leading pressure ($\VPij[3]=-B_{ij}$) \\
$\VPikj$ & $\int_0^1\Theta_{kj}\,\varphi_i\,d\sigma$ & $3$ & Leading pressure, non-linear \\
\midrule
\multicolumn{4}{l}{\textit{Static horizontal tensors (assembled once)}}\\
$\HMsm$ & $\int_{\Omega_h}\psi_s\psi_m\,d\Omega$ & $2$ & Horizontal mass (continuity) \\
$\HMsnm$ & $\int_{\Omega_h}\psi_s\psi_n\psi_m\,d\Omega$ & $3$ & $H$-weighted mass carrier \\
$\HDsnma$ & $\int_{\Omega_h}\psi_m\psi_n\nablah\psi_s\,d\Omega$ & $3{+}$ & Continuity flux (test grad.) \\
$\mathbb{E}^H_{sm\gamma}$ & $\int_{\Omega_h}\psi_s\psi_m\nablah\psi_\gamma\,d\Omega$ & $3{+}$ & Gravity (trial grad.) \\
$\mathbb{B}^h_{sm}$ & $\int_{\Omega_h}\psi_s\psi_m\nablah h\,d\Omega$ & $2{+}$ & Bed-slope carrier \\
$\mathbb{K}^H_{sn}$ & $\int_{\Omega_h}\nablah\psi_s\otimes\nablah\psi_n\,d\Omega$ & $2{+}$ & Stiffness (flat-bed dispersion) \\
$\mathbb{T}^H_{smnr}$ & $\int_{\Omega_h}\psi_s\psi_m\psi_n\nablah\psi_r\,d\Omega$ & $4{+}$ & Advection (formal; not stored) \\
\midrule
\multicolumn{4}{l}{\textit{State-weighted horizontal operators (re-assembled per step / iterate)}}\\
$\HmMdwH$ & $\int_{\Omega_h}H\,\psi_s\psi_n\,d\Omega$ & $2$ & Acceleration mass \\
$\HDdwsn(H)$ & $\int_{\Omega_h}H\,\psi_n\nablah\psi_s\,d\Omega$ & $2{+}$ & Continuity flux \\
$\bar{\mathbf{C}}^H_{sr}(H,\uk)$ & $\int_{\Omega_h}\psi_s\,H(\uk\cdot\nablah\psi_r)\,d\Omega$ & $2$ & Horizontal advection \\
$\bar{\mathbf{S}}^H_{sr}(H,\uk)$ & $\int_{\Omega_h}\psi_s\,(\nablah\cdot[H\uk])\,\psi_r\,d\Omega$ & $2$ & Vertical ($\omega$) advection \\
$\bar{\mathbf{P}}^H_{sn}(H)$ & $\int_{\Omega_h}H^2\nablah\psi_s\otimes(H\nablah\psi_n+\psi_n\nablah H)\,d\Omega$ & $2{+}$ & Leading pressure (dispersion) \\
$\bar{\mathbb{A}}^{(\ell)}_{sn},\,\bar{\mathbb{K}}^{(\ell)}_{sn}$ & eq.~\eqref{eq: def state weighted linear pressure operators} & $2{+}$ & Linear slope pressure \\
\bottomrule
\end{tabular}
\caption{Complete list of pre-computable tensors and per-step operators for the LFE-M model. Vertical tensors depend only on the vertical mesh $\{c_k\}$ and shape functions $\phi_j(\sigma)$ and are assembled once. Static horizontal tensors depend only on the horizontal mesh and basis $\psi_n(\mathbf{x})$ and are assembled once; state-weighted operators additionally carry the current state and are re-assembled cheaply each step. A ``$+$'' marks operators with a free spatial component index (vector-/matrix-valued integrals).}
\label{tab: all tensors}
\end{table}


\end{document}
